\subsection{Ranks of projective modules}

DEFINITION 1. Let P be a finitely generated projective A-module. For every prime ideal p of A, the rank of the free \(A_{p}\) -module \(P_{p}\) is called the rank of P at p and is denoted by \(\operatorname{rg}_{\mathfrak{p}}(\mathbf{P})\) .

By Theorem 1 the integer-valued function \(\mathfrak{p} \mapsto \mathrm{rg}_{\mathfrak{p}}(\mathrm{P})\) is locally constant on \(\mathbf{X} = \operatorname{Spec}(\mathbf{A})\) ; it is therefore constant if \(\mathbf{X}\) is connected and in particular if the ring \(\mathbf{A}\) is an integral domain (§ 4, no. 3, Corollary 2 to Proposition 15).

DEFINITION 2. Let \(n\) be an integer \(\geqslant 0\) . A projective A-module \(P\) is said to be of rank \(n\) if it is finitely generated and \(\operatorname{rg}_{\mathfrak{p}}(P) = n\) for every prime ideal \(\mathfrak{p}\) of \(A\) .

Clearly every finitely generated free A-module L is of rank n in the sense of Definition 2, n being equal to the dimension (or rank) of L defined in Algebra, Chapter II, §7, no. 2.

A projective module of rank 0 is zero (§ 3, no. 3, Corollary 2 to Theorem 1). If A is not reduced to 0 and a projective A-module P is of rank n, the integer n is determined uniquely; it is then denoted by rg(P).

THEOREM 2. Let P be an A-module and n an integer \(\geqslant 0\) . The following properties are equivalent:

\begin{enumerate}
\def\labelenumi{(\alph{enumi})}
\item
  P is projective of rank n.
\item
  \(\mathbf{P}\) is finitely generated and, for every maximal ideal \(\mathfrak{m}\) of \(\mathbf{A}\) , the \(\mathbf{A}_{\mathfrak{m}}\) -module \(\mathbf{P}_{\mathfrak{m}}\) is free of rank \(n\) .
\item
  \(\mathbf{P}\) is finitely generated and, for every prime ideal \(\mathfrak{p}\) of \(\mathbf{A}\) , the \(\mathbf{A}_{\mathfrak{p}}\) -module \(\mathbf{P}_{\mathfrak{p}}\) is free of rank \(n\) .
\item
  For every maximal ideal m of A, there exists \(f \in A - m\) such that the \(A_{f}\) -module \(P_{f}\) is free of rank n.
\end{enumerate}

By Definition 2 and Theorem 1, (a) and (c) are equivalent; (b) implies (c), as, for every prime ideal p of A, there exists a maximal ideal m containing p and, writing \(p' = p_{m}\) , \(P_{p}\) is isomorphic to \((\mathrm{P}_{\mathfrak{m}})_{\mathfrak{p'}}\) (§ 2, no. 5, Proposition 11); if \(P_{m}\) is free of rank n, so then is \(P_{p}\) . Property (c) implies (d) by virtue of Theorem 1 and the fact that, if \(f \in A - m\) and \(m' = m_{f}\) , \(P_{m}\) is isomorphic to \((\mathrm{P}_{f})_{\mathfrak{m'}}\) and therefore the ranks of \(P_{f}\) and \(P_{m}\) are equal. Finally, this last argument and Theorem 1 show that (d) implies (b).

Remark. If A is an integral domain, a projective A-module admits a well-defined rank (in the sense of Definition 2), as has been observed above; moreover, this rank coincides with the rank defined in Algebra, Chapter II, §7, no. 2; it is sufficient to apply Theorem 2 (c) with p = (0).

Let E and F be two finitely generated projective A-modules. We know (Algebra, Chapter II, §§ 2 and 3) that \(\mathbf{E} \times \mathbf{F}\) , \(\mathbf{E} \otimes_{\mathbf{A}} \mathbf{F}\) , \(\operatorname{Hom}_{\mathbf{A}}(\mathbf{E}, \mathbf{F})\) and the dual \(\mathbf{E}^*\)

of E are projective and finitely generated; so is the exterior power \(\wedge^k\) E for every integer \(k > 0\) (Algebra, Chapter III). Also, it follows immediately from Definition 1 and §2, no. 7, Propositions 18 and 19 and no. 8, that, for every prime ideal \(\mathfrak{p}\) of A:

\[
\mathrm{rg} _ {\mathfrak {p}} (\mathrm{E} \times \mathrm{F}) = \mathrm{rg} _ {\mathfrak {p}} (\mathrm{E}) + \mathrm{rg} _ {\mathfrak {p}} (\mathrm{F})\tag{1}
\]

\[
\operatorname{rg} _ {\mathfrak {p}} (\mathrm{E} \otimes_ {\mathrm{A}} \mathrm{F}) = \operatorname{rg} _ {\mathfrak {p}} (\mathrm{E}). \operatorname{rg} _ {\mathfrak {p}} (\mathrm{F})\tag{2}
\]

\[
\operatorname{rg} _ {\mathfrak {p}} (\operatorname{Hom} _ {\mathbf {A}} (\mathrm{E}, \mathrm{F})) = \operatorname{rg} _ {\mathfrak {p}} (\mathrm{E}). \operatorname{rg} _ {\mathfrak {p}} (\mathrm{F})\tag{3}
\]

\[
\operatorname{rg} _ {\mathfrak {p}} (\mathrm{E} ^ {*}) = \operatorname{rg} _ {\mathfrak {p}} (\mathrm{E})\tag{4}
\]

\[
\operatorname{rg} _ {\mathfrak {p}} (\overset {k} {\wedge} \mathrm{E}) = \binom{\operatorname{rg} _ {\mathfrak {p}} (\mathrm{E})}{k}.\tag{5}
\]

If the ranks of E and F are defined, so are those of \(E \times F\) , \(E \otimes_{A} F\) , \(\operatorname{Hom}_{A}(E, F)\) , \(E^{*}\) and \(\Lambda^{k} E\) and the above equations also hold with the index p omitted. Moreover:

COROLLARY. For a finitely generated projective A-module P to be of rank n, it is necessary and sufficient that \(\wedge^{n}P\) be of rank 1.

PROPOSITION 4. Let B be a commutative A-algebra and P a projective A-module of rank n.~The B-module \(\mathbf{P}_{(\mathbf{B})} = \mathbf{B} \otimes_{\mathbf{A}} \mathbf{P}\) is then projective of rank n.

We know that \(P_{(B)}\) is projective and finitely generated (Algebra, Chapter II, §5, no. 1, Corollary to Proposition 4). If q is a prime ideal of B and p its inverse image in A, then

\[
\left(\mathrm{P} _ {(\mathrm{B})}\right) _ {\mathrm{q}} = \left(\mathrm{P} \otimes_ {\mathrm{A}} \mathrm{B}\right) \otimes_ {\mathrm{B}} \mathrm{B} _ {\mathrm{q}} = \mathrm{P} \otimes_ {\mathrm{A}} \mathrm{B} _ {\mathrm{q}} = \left(\mathrm{P} \otimes_ {\mathrm{A}} \mathrm{A} _ {\mathrm{p}}\right) \otimes_ {\mathrm{A}} \mathrm{B} _ {\mathrm{q}}
\]

and, as, by hypothesis, \(\mathbf{P} \otimes_{\mathbf{A}} \mathbf{A}_{\mathfrak{p}}\) is a free \(\mathbf{A}_{\mathfrak{p}}\) -module of rank \(n\) , \((\mathrm{P}_{(\mathbf{B})})_{\mathfrak{q}}\) is a free \(\mathbf{B}_{\mathfrak{q}}\) -module of rank \(n\) .

PROPOSITION 5. Let A be a semi-local ring and P a finitely generated projective A-module. If the rank of P is defined, P is a free A-module.

Suppose first that A is isomorphic to a product of fields \(K_{i}\) ( \(1 \leqslant i \leqslant n\) ). The \(K_{i}\) are then identified with the minimal ideals (Algebra, Chapter VIII, §3, no. 1) of A and, for all i, the sum \(p_{i}\) of the \(K_{j}\) of index \(j \neq i\) is a maximal ideal of A, the \(p_{i}\) ( \(1 \leqslant i \leqslant n\) ) being the only prime ideals of A. Every finitely generated A-module P is therefore the direct sum of its isotypical components \(P_{i}\) ( \(1 \leqslant i \leqslant n\) ), \(P_{i}\) being isomorphic to a direct sum of a finite number \(r_{i}\) of A-modules isomorphic to \(K_{i}\) (Algebra, Chapter VIII, §5, no. 1, Proposition 1 and no. 3, Proposition 11); the ring \(A_{p_{i}}\) is identified with \(K_{i}\) and annihilates the \(P_{j}\) of index \(j \neq i\) , hence \(r_{i} = \operatorname{rg}_{\mathfrak{p}_{i}}(\mathbf{P})\) ; if all the \(r_{i}\) are equal to the same number r, P is isomorphic to \(A^{r}\) , whence the proposition in this case. In the general case, let R be the Jacobson radical of A and B = A/R; as B is a product of fields, the projective B-module \(P_{(B)}\) is free by the remarks preceding Proposition 4. Also P is a flat A-module and the proposition then follows from §3, no. 2, Proposition 5.

\subsection{Projective modules of rank 1}

THEOREM 3. Let A be a ring and M a finitely generated A-module.

\begin{enumerate}
\def\labelenumi{(\roman{enumi})}
\item
  If there exists an A-module N such that \(\mathbf{M} \otimes_{\mathbf{A}} \mathbf{N}\) is isomorphic to A, the module M is projective of rank 1.
\item
  Conversely, if \(\mathbf{M}\) is projective of rank 1 and \(\mathbf{M}^*\) is the dual of \(\mathbf{M}\) , the canonical homomorphism \(u: \mathbf{M} \otimes_{\mathbf{A}} \mathbf{M}^* \to \mathbf{A}\) corresponding to the canonical bilinear form \((x, x^*) \to \langle x, x^* \rangle\) on \(\mathbf{M} \times \mathbf{M}^*\) (Algebra, Chapter II, §2, no. 3) is bijective.
\item
  It is required to prove that, for every maximal ideal m of A, the \(A_{m}\) -module \(M_{m}\) is free of rank 1 (Theorem 2 (b)); we are free to replace A by \(A_{m}\) and hence may assume that A is a local ring ( \(\S\) 2, no. 7, Proposition 18). Let k = A/m. The isomorphism v: \(M \otimes_{A} N \to A\) defines an isomorphism
\end{enumerate}

\[
v \otimes 1 _ {k} \colon (\mathbf {M} / \mathfrak {m M}) \otimes_ {k} (\mathbf {N} / \mathfrak {m N}) \rightarrow k
\]

as the rank over \(k\) of (M/mM) \(\otimes_k\) (N/mN) is the product of the ranks of M/mM and N/mN, these latter are necessarily equal to 1, in other words M/mM is monogenous. It follows that M is monogenous (§ 3, no. 2, Corollary 2 to Proposition 4); on the other hand, the annihilator of M also annihilates M \(\otimes_A\) N and hence is zero, which proves that M is isomorphic to A.

\begin{enumerate}
\def\labelenumi{(\roman{enumi})}
\setcounter{enumi}{1}
\tightlist
\item
  It is sufficient to prove that, for every maximal ideal m of A, \(u_{m}\) is an isomorphism ( \(\S 3\) , no. 3, Theorem 1). As M is finitely presented (Chapter I, \(\S 2\) , no. 8, Lemma 8), \((\mathbf{M}^{*})_{\mathfrak{m}}\) is canonically identified with the dual \((\mathbf{M}_{\mathfrak{m}})^{*}\) ( \(\S 2\) , no. 7, Proposition 19) and, as \(M_{m}\) is free of rank 1 like its dual \((\mathbf{M}_{\mathfrak{m}})^{*}\) , clearly the canonical homomorphism \(u_{\mathfrak{m}}: (\mathbf{M}_{\mathfrak{m}}) \otimes_{\mathbf{A}\mathfrak{m}} (\mathbf{M}_{\mathfrak{m}})^{*} \to \mathbf{A}_{\mathfrak{m}}\) is bijective, which completes the proof.
\end{enumerate}

Remark (1). If \(M\) is projective of rank 1 and \(\mathbf{N}\) is such that \(\mathbf{M} \otimes_{\mathbf{A}} \mathbf{N}\) is isomorphic to \(\mathbf{A}\) , then \(\mathbf{N}\) is isomorphic to \(\mathbf{M}^*\) : there are isomorphisms

\[
\mathbf {N} \rightarrow \mathbf {N} \otimes \mathbf {A} \rightarrow \mathbf {N} \otimes \mathbf {M} \otimes \mathbf {M} ^ {*} \rightarrow \mathbf {A} \otimes \mathbf {M} ^ {*} \rightarrow \mathbf {M} ^ {*}.
\]

PROPOSITION 6. Let M and N be projective A-modules of rank 1. Then M ⊗\_\{A\} N, Hom\_\{A\}(M, N) and the dual M* of M are projective of rank 1.

This follows immediately from formulae (2), (3) and (4).

Let us now note that every finitely generated A-module is isomorphic to a quotient module of \(L = A^{(N)}\) ; we may therefore speak of the set \(F(A)\) of classes of finitely generated A-modules with respect to isomorphism (Set Theory, Chapter I, §6, no. 9); we denote by \(P(A)\) the subset of \(F(A)\) consisting of the classes of projective A-modules of rank 1 and by \(\operatorname{cl}(M)\) the image in \(P(A)\) of a projective A-module M of rank 1. It is immediate that, for two projective A-modules M,

N of rank 1, \(\operatorname{cl}(\mathbf{M} \otimes_{\mathbf{A}} \mathbf{N})\) depends only on \(\operatorname{cl}(\mathbf{M})\) and \(\operatorname{cl}(\mathbf{N})\) ; as definition we set

\[
\operatorname{cl} (\mathbf {M}) + \operatorname{cl} (\mathbf {N}) = \operatorname{cl} (\mathbf {M} \otimes_ {\mathbf {A}} \mathbf {N})\tag{6}
\]

and an internal law of composition is thus defined on \(\pmb{P}(\mathbf{A})\) .

PROPOSITION 7. The set P(A) of classes of projective A-modules of rank 1, with the law of composition (6), is a commutative group. If M is a projective A-module of rank 1 and M* is its dual, then

\[
\operatorname{cl} (\mathbf {M} ^ {*}) = - \operatorname{cl} (\mathbf {M}) \quad a n d \quad \operatorname{cl} (\mathbf {A}) = 0.\tag{7}
\]

The associativity and commutativity of the tensor product show that the law of composition (6) is associative and commutative; the isomorphism between \(A \otimes_{A} M\) and M prove that \(\operatorname{cl}(A)\) is the identity element under this law and, by virtue of Theorem 3, \(\operatorname{cl}(M) + \operatorname{cl}(M^{*}) = \operatorname{cl}(A)\) , whence the proposition.

Let B be a commutative A-algebra and M a projective A-module of rank 1; then \(\mathbf{M}_{(\mathbf{B})} = \mathbf{B} \otimes_{\mathbf{A}} \mathbf{M}\) is a projective B-module of rank 1 (no. 3, Proposition 4). Then there exists a mapping called canonical \(\phi: P(\mathbf{A}) \to P(\mathbf{B})\) such that

\[
\phi (\operatorname{cl} (M)) = \operatorname{cl} (M _ {(B)}).\tag{8}
\]

The equation \(\mathbf{M}_{(\mathbf{B})}\otimes_{\mathbf{B}}\mathbf{N}_{(\mathbf{B})}=(\mathbf{M}\otimes_{\mathbf{A}}\mathbf{N})_{(\mathbf{B})}\) for two A-modules M, N proves that the mapping \(\phi\) is a commutative group homomorphism.

*Remark (2). Condition (e) of Theorem 1 (equivalent to the fact that P is projective and finitely generated) may also be expressed by saying that the sheaf of modules \(\tilde{\mathbf{P}}\) over \(\mathbf{X} = \operatorname{Spec}(\mathbf{A})\) associated (*) with P is locally free and of finite type and may therefore be interpreted as the sheaf of sections of a vector bundle over X. Conversely, every vector bundle over X arises from a finitely generated projective module, which is determined to within a unique isomorphism; the projective modules of rank \(n\) thus correspond to the vector bundles all of whose fibres have dimension \(n\) . In particular, the vector bundles of rank 1 correspond to the projective modules of rank 1. If we denote by \(\mathcal{O}_{\mathbf{x}}\) the structure sheaf \(\tilde{\mathbf{A}}\) and by \(\mathcal{O}_{\mathbf{x}}^{*}\) the sheaf of units of \(\mathcal{O}_{\mathbf{x}}\) (whose sections over an open set U of X are the invertible elements of the ring of sections of \(\mathcal{O}_{\mathbf{x}}\) over U), it follows that the group \(P(\mathbf{A})\) is isomorphic to the first cohomology group \(\mathrm{H}^1 (\mathrm{X},\mathcal{O}_\mathrm{x}^*)\) .

\subsection{Non-degenerate submodules}

In this no. and the two following, A denotes a ring, S a multiplicative subset of A consisting of elements which are not divisors of zero in A, and B the ring \(S^{-1}A\) ; A is canonically identified with a subring of B ( \(\S\) 2, no. 1, Remark 3). The elements of S are therefore invertible in B.

One of the most important special cases for applications is that where A is an integral domain and S is the set of elements \(\neq0\) of A; B is then the field of fractions of A.

DEFINITION 3. Let M be a sub-A-module of B. M is called non-degenerate if B. M = B.

If B is a field, this condition simply means that M is not reduced to 0.

PROPOSITION 8. Let M be a sub-A-module of B. The following conditions are equivalent:

\begin{enumerate}
\def\labelenumi{(\alph{enumi})}
\item
  M is non-degenerate.
\item
  M meets S.
\item
  If \(j: \mathbf{M} \to \mathbf{B}\) is the canonical injection, the homomorphism \(u = \mathbf{S}^{-1}j: \mathbf{S}^{-1}\mathbf{M} \to \mathbf{B}\) is bijective.
  (a) implies (b), for if B.M = B, there exists \(a \in A\) , \(s \in S\) and \(x \in M\) such that \((a/s)x = 1\) , hence ax = s belongs to S ∩ M. To see that (b) implies (c), note that u is already injective (§2, no. 4, Theorem 1); moreover, if \(x \in M \cap S\) , the image under u of \(x/x \in S^{-1}M\) in B is equal to 1 and u is therefore surjective. Finally, (c) clearly implies (a).
\end{enumerate}

COROLLARY. If M and N are two non-degenerate sub-A-modules of B, the A-modules M + N, M.N and M ∩ N are non-degenerate.

The assertion is trivial for \(\mathbf{M} + \mathbf{N}\) ; on the other hand if \(s \in \mathbf{S} \cap \mathbf{M}\) and \(t \in \mathbf{S} \cap \mathbf{N}\) , then \(st \in \mathbf{S} \cap (\mathbf{M}. \mathbf{N})\) and \(st \in \mathbf{S} \cap (\mathbf{M} \cap \mathbf{N})\) .

Given two sub-A-modules M and N of B, let us denote by N: M the sub-A-module of B consisting of those \(b \in B\) such that \(bM \subset N\) (Chapter I, §2, no. 10, Remark). If every \(b \in N\) : M is mapped to the homomorphism \(h_{b}: x \mapsto bx\) of M to N, a canonical homomorphism \(b \mapsto h_{b}\) is obtained from N: M to \(\operatorname{Hom}_{\mathbf{A}}(\mathbf{M}, \mathbf{N})\) .

PROPOSITION 9. Let M, N be two sub-A-modules of B. If M is non-degenerate, the canonical homomorphism from N: M to \(\mathrm{Hom}_{\mathbf{A}}(\mathbf{M},\mathbf{N})\) is bijective.

Let \(s \in S \cap M\) . If \(b \in N: M\) is such that bx = 0 for all \(x \in M\) , then bs = 0, whence b = 0 since s is not a divisor of 0 in B. On the other hand, let \(f \in \operatorname{Hom}_{\mathbf{A}}(\mathbf{M}, \mathbf{N})\) and set \(b = f(s) / s\) ; for all \(x \in \mathbf{M}\) , there exists \(t \in \mathbf{S}\) such that \(tx \in \mathbf{A}\) . Then

\[
f (x) = s ^ {- 1} t ^ {- 1} f (s t x) = s ^ {- 1} t ^ {- 1} t x f (s) = b x,
\]

whence \(b \in \mathbf{N}\) : M and \(f = h_b\) , which proves the proposition.

Remark. In particular, A: M is canonically identified with the dual M* of M, the canonical bilinear form on M × M* being identified with the restriction to M × (A: M) of the multiplication B × B → B.

\subsection{Invertible submodules}

*(We preserve the notation of no. 5.)*

DEFINITION 4. A sub-A-module M of B is called invertible if there exists a sub-A-module N of B such that M.N = A.

Example. If \(b\) is invertible element of \(\mathbf{B}\) , the A-module \(Ab\) is invertible, as is seen by taking \(\mathbf{N} = \mathbf{Ab}^{-1}\) .

PROPOSITION 10. Let M be an invertible sub-A-module of B. Then:

\begin{enumerate}
\def\labelenumi{(\roman{enumi})}
\item
  There exists \(s \in S\) such that \(As \subset M \subset As^{-1}\) (and in particular \(M\) is nondegenerate).
\item
  A: M is the only sub-A-module N of B such that M.N = A.
\end{enumerate}

If M.N = A, then B.M = B.(B.M) ⊃ B.(M.N) = B.A = B, hence M is non-degenerate. Similarly N is non-degenerate. If \(t \in S \cap M\) and \(u \in S \cap N\) (no. 5, Proposition 8), the element s = tu belongs to \(S \cap M \cap N\) , whence \(Ms \subset M.N = A\) and therefore \(As \subset M \subset As^{-1}\) .

On the other hand, obviously \(N \subset A: M\) , whence

\[
\mathbf {A} = \mathbf {M}. \mathbf {N} \subset \mathbf {M}. (\mathbf {A}: \mathbf {M}) \subset \mathbf {A}
\]

and M. (A: M) = A; multiplying the two sides by N, we deduce A: M = N, which completes the proof.

THEOREM 4. Let M be a non-degenerate sub-A-module of B. The following properties are equivalent:

\begin{enumerate}
\def\labelenumi{(\alph{enumi})}
\item
  M is invertible.
\item
  M is projective.
\item
  M is projective of rank 1.
\item
  M is a finitely generated A-module and, for every maximal ideal m of A, the \(A_{m}\) -module \(M_{m}\) is monogenous.
\end{enumerate}

Let us show first the equivalence of properties (a), (b) and (c). If (a) holds and N is sub-A-module of B such that M.N = A, then there is a relation

\[
\sum_ {i = 1} ^ {p} m _ {i} n _ {i} = 1 \quad (m _ {i} \in \mathbf {M}, n _ {i} \in \mathbf {N} \text {   for   all   } i).\tag{9}
\]

\[
\boldsymbol {x} \in \mathbf {M},
\]

\[
v _ {i} (x) = n _ {i} x;
\]

\[
\nu_ {i}
\]

\(x = \sum_{i=1}^{n} m_i v_i(x)\) for all \(x \in \mathbf{M}\) ; this proves (Algebra, Chapter II, § 2, no. 6, Proposition 12) that \(\mathbf{M}\) is projective and generated by the \(m_i\) ; hence \(\mathbf{M}\) is a finitely generated projective module.

Let \(\mathfrak{m}\) be a maximal ideal of \(\mathbf{A}\) ; we show that the integer \(r = \operatorname{rg}_{\mathfrak{m}}(\mathbf{M})\) is equal to 1. Let \(S'\) be the image of \(S\) in \(A_{\mathfrak{m}}\) ; as the elements of \(S\) are not divisors of 0 in \(A\) , those of \(S'\) are not divisors of 0 in \(A_{\mathfrak{m}}\) , since \(A_{\mathfrak{m}}\) is a flat \(A\) -module (§2, no. 4, Theorem 1 and Chapter I, §2, no. 4, Proposition 3); then \(S'^{-1}A_{\mathfrak{m}} \neq 0\) and, as \(M_{\mathfrak{m}}\) is a free \(A_{\mathfrak{m}}\) -module of rank \(r\) , \(S'^{-1}M_{\mathfrak{m}}\) is a free \(S'^{-1}A_{\mathfrak{m}}\) -module of rank \(r\) . But if \(T'\) is the image of \(A - \mathfrak{m}\) in \(S^{-1}A\) , \(S'^{-1}A_{\mathfrak{m}}\) (resp. \(S'^{-1}M_{\mathfrak{m}}\) ) is canonically identified with \(T'^{-1}(S^{-1}A)\) (resp. \(T'^{-1}(S^{-1}M)\) ) (§2, no. 3, Proposition 7). Now \(S^{-1}M = B\) (Proposition 8 (c)) and hence \(T'^{-1}(S^{-1}M)\) is a free \(A\) -module of rank 1 over \(T'^{-1}(S^{-1}A)\) , which proves that \(r = 1\) and shows the implication (a) ⇒ (c).

The implication (c) \(\Rightarrow\) (b) is trivial. Let us show that (b) \(\Rightarrow\) (a). There exists by hypothesis a family (not necessarily finite) \((f_{\lambda})_{\lambda \in L}\) of linear forms on M and a family \((m_{\lambda})_{\lambda \in L}\) of elements of M such that, for all \(x \in M\) , the family \((f_{\lambda}(x))\) has finite support and \(x = \sum_{\lambda \in L} m_{\lambda} f_{\lambda}(x)\) (Algebra, Chapter II, §2, no. 6, Proposition 12). Since M is non-degenerate, \(f_{\lambda}(x) = n_{\lambda} x\) for some \(n_{\lambda} \in A: M\) by virtue of Proposition 9 of no. 5. Taking \(x\) as an element of M \(\cap S\) (no. 5, Proposition 8), it is seen that of necessity \(n_{\lambda} = 0\) except for a finite number of indices and \(\sum_{\lambda \in L} m_{\lambda} n_{\lambda} = 1\) . This obviously implies M. (A: M) = A, whence (a).

By virtue of Definition 2 of no. 3, (c) implies (d). Let us show the converse. As M is non-degenerate, its annihilator is zero (Proposition 8 (b)), then so is the annihilator of \(M_{m}\) (§ 2, no. 4, formula (9)). As \(M_{m}\) is assumed to be a monogenous \(A_{m}\) -module, it is therefore free of rank 1 and it then follows from no. 3, Theorem 2 that M is projective of rank 1.

COROLLARY. Every invertible sub-A-module of B is flat and finitely presented. This follows from Theorem 4 (c).

PROPOSITION 11. Let M, N be two sub-A-modules of B. Suppose that M is invertible. Then:

\begin{enumerate}
\def\labelenumi{(\roman{enumi})}
\tightlist
\item
  The canonical homomorphism \(\mathbf{M} \otimes_{\mathbf{A}} \mathbf{N} \to \mathbf{M}. \mathbf{N}\) is bijective.
\end{enumerate}

\[
\text {(ii)} \mathbf {N}: \mathbf {M} = \mathbf {N}. (\mathbf {A}: \mathbf {M}) a n d \mathbf {N} = (\mathbf {N}: \mathbf {M}). \mathbf {M}.
\]

Let \(j\) be the canonical injection \(N \to B\) . Since \(M\) is a flat A-module (Corollary to Theorem 4), \(1 \otimes j \colon M \otimes_{A} N \to M \otimes_{A} B\) is injective. But, as \(B = S^{-1}A\) , the B-module \(M \otimes_{A} B\) is equal to \(S^{-1}M\) and hence is identified with \(B\) since \(M\) is non-degenerate (no. 5, Proposition 8). If this identification is made, the image of \(1 \otimes j\) is \(M \cdot N\) , whence (i).

Let us set \(\mathbf{M}' = \mathbf{A} : \mathbf{M}\) . Then obviously \(\mathbf{M}' . \mathbf{N} \subset \mathbf{N} : \mathbf{M}\) and \(\mathbf{M} . (\mathbf{N} : \mathbf{M}) \subset \mathbf{N}\) . On the other hand, since \(\mathbf{M} . \mathbf{M}' = \mathbf{A}\) (Proposition 10),

\[
\mathbf {N}: \mathbf {M} = \mathbf {M} ^ {\prime}. \mathbf {M}. (\mathbf {N}: \mathbf {M}) \subset \mathbf {M} ^ {\prime}. \mathbf {N}
\]

and \(\mathbf{N} = \mathbf{M}.\mathbf{M}'.\mathbf{N}\subset \mathbf{M}.(\mathbf{N}:\mathbf{M})\) , whence (ii).

Remark. The proof of (i) in Proposition 11 uses only the fact that \(\mathbf{M}\) is flat and non-degenerate.

\subsection{The group of classes of invertible modules}

*(We preserve the notation of nos. 5 and 6.)*

Under multiplication, the sub-A-modules of B form a commutative monoid M, with A as identity element. Then the invertible modules are the invertible elements of M and therefore form a commutative group J. We have seen (no. 6, Proposition 10) that the inverse of \(M \in J\) is A: M.

Let A* (resp. B*) be the multiplicative group of invertible elements of A (resp. B) and let u denote the canonical injection A → B. For all \(b \in B^{*}\) , \(\theta(b) = bA\) is an invertible sub-A-module. The mapping \(\theta: B^{*} \to J\) is a homomorphism whose kernel is \(u(A^{*})\) ; its cokernel will be denoted by C or C(A). The group C is called the group of classes of invertible sub-A-modules of B. The following exact sequence has been constructed

\[
(1) \longrightarrow A ^ {*} \stackrel {u} {\longrightarrow} B ^ {*} \stackrel {\theta} {\longrightarrow} \mathfrak {J} \stackrel {\rho} {\longrightarrow} \mathfrak {C} \longrightarrow (1)\tag{10}
\]

where (1) denotes the group consisting only of the identity element and \(\rho\) is the canonical mapping \(\mathfrak{J} \to \mathfrak{C} = \mathfrak{J}/\theta(\mathrm{B}^*)\) .

As every invertible sub-A-module M of B is projective of rank 1 (no. 6, Theorem 4), the element \(\operatorname{cl}(M) \in P(A)\) is defined (no. 4).

PROPOSITION 12. The mapping cl: \(\mathfrak{J} \to P(A)\) defines, by taking quotients, an isomorphism from \(\mathfrak{C} = \mathfrak{J}/\theta(B^{*})\) onto the kernel of the canonical homomorphism

\[
\phi \colon P (\mathrm{A}) \to P (\mathrm{B})
\]

(no. 4).

In other words, there is an exact sequence

\[
(1) \longrightarrow \mathrm{A} ^ {*} \stackrel {u} {\longrightarrow} \mathrm{B} ^ {*} \stackrel {\theta} {\longrightarrow} \mathfrak {J} \stackrel {\mathrm{cl}} {\longrightarrow} P (\mathrm{A}) \stackrel {\phi} {\longrightarrow} P (\mathrm{B}).\tag{11}
\]

It follows from Proposition 11 of no. 6 and the definition of addition in \(\mathbf{P}(\mathbf{A})\) that \(\operatorname{cl}(\mathbf{M},\mathbf{N}) = \operatorname{cl}(\mathbf{M}) + \operatorname{cl}(\mathbf{N})\) for \(\mathbf{M}\) , \(\mathbf{N}\) in \(\mathfrak{J}\) , which shows that \(\operatorname{cl}\) is a homomorphism. If \(\mathbf{M} \in \mathfrak{J}\) is isomorphic to \(\mathbf{A}\) , there exists \(b \in \mathbf{B}\) such that \(\mathbf{M} = Ab\) and, as \(\mathbf{M}\) is invertible, there exists \(b' \in B\) such that \(b'b = 1\) , in other words \(b\) is invertible in \(\mathbf{B}\) ; the converse is immediate. Hence the kernel of \(\operatorname{cl}\) in \(\mathfrak{J}\) is \(\theta(\mathbf{B}^*)\) .

Let us now determine the image of cl. If \(M \in J\) , then \(M \otimes_{A} B = S^{-1}M = B\) (no. 5, Proposition 8 (c)), whence \(\operatorname{cl}(M) \in \operatorname{Ker}(\phi)\) . Conversely, let P be a projective A-module of rank 1 such that \(P_{(B)} = P \otimes_{A} B\) is B-isomorphic to B. As P is a flat A-module, the injection \(u: A \to B\) defines an injection \(u \otimes 1: P \to P_{(B)} = B\) and P is thus identified with a sub-A-module of B; by virtue of Proposition 8 (c) of no. 5, P is non-degenerate and Theorem 4 of no. 6 shows that P is invertible. The kernel of \(\phi\) is therefore equal to the image of \(\operatorname{cl}: J \to P(A)\) .

COROLLARY 1. For two invertible sub-A-modules of B to have the same image in C, it is necessary and sufficient that they be isomorphic.

COROLLARY 2. If the ring B is semi-local, the group C of classes of invertible sub-A-modules of B is canonically identified with the group P(A) of classes of projective A-modules of rank 1.

In this case \(\mathrm{P(B)} = 0\) (no. 3, Proposition 5).

Remark. The hypothesis of Corollary 2 is fulfilled in the two following cases:

\begin{enumerate}
\def\labelenumi{(\arabic{enumi})}
\tightlist
\item
  A is an integral domain and S is the set of elements \(\neq0\) of A, B then being the field of fractions of A. The invertible sub-A-modules of B are also called in this case invertible fractional ideals; those which are monogenous free A-modules Ab ( \(b\neq0\) in B) are just the fractional principal ideals defined in Algebra, Chapter VI, §1, no. 5.
\end{enumerate}

*(2) The ring A is Noetherian and S is the set of elements of A which are not divisors of 0 such that B is the total ring of fractions of A. In this case \(S = A - \bigcup_{i} p_i\) , where the \(p_i\) are the elements (finite in number) of Ass(A) (Chapter IV, §1), hence B is semi-local (§3, no. 5, Proposition 17). *


\exercisesection{1}

¶ 1. (a) Show that a group G cannot be the union of two subgroups distinct from G. Show that, for every set I with at least two elements, the commutative group \(G = \mathbf{F}_{2}^{(I)}\) is the union of three subgroups distinct from G.

\begin{enumerate}
\def\labelenumi{(\alph{enumi})}
\setcounter{enumi}{1}
\item
  Let \((\mathbf{H}_i)_{i\in I}\) be a finite family of subgroups of a group \(G\) such that each of the \(\mathbf{H}_i\) is a subgroup of \(G\) of infinite index. Show that \(G\) cannot be the union of a finite number of left cosets of the \(\mathbf{H}_i\) . (Argue by induction on the number of elements in \(I\) ; if there exist two distinct indices \(i,j\) such that the index \((\mathbf{H}_i:(\mathbf{H}_i\cap \mathbf{H}_j))\) is finite, \(\mathbf{H}_i\) may be suppressed; if on the other hand \(\mathbf{H}_i\cap \mathbf{H}_j\) is of infinite index in \(\mathbf{H}_i\) for every ordered pair \((i,j)\) of distinct indices, consider an index \(k\) such that \(\mathbf{H}_k\) is maximal in the set of \(\mathbf{H}_i\) and show that \(\mathbf{H}_k\) is the union of a finite number of left cosets of the \(\mathbf{H}_k\cap \mathbf{H}_i\) where \(i\neq k\) .)
\item
  Give an example of a commutative ring A and four ideals \(a, b_{1}, b_{2}, b_{3}\) of A such that \(a \notin b_{i} (i = 1, 2, 3)\) , but \(a = \bigcup_{i} b_{i}\) (use (a)).
\end{enumerate}

\begin{enumerate}
\def\labelenumi{\arabic{enumi}.}
\setcounter{enumi}{1}
\item
  Let A be a ring which is not necessarily commutative and \(a, p_{1}, \ldots, p_{n}\) two-sided ideals of A. Suppose that a is contained in the union of the \(p_{i}\) and that all the \(p_{i}\) except at most two are prime ideals (Algebra, Chapter VIII, §8, Exercise 6). Show that a is contained in one of the \(p_{i}\) .
\item
  In the product ring \(\mathbf{A} = \mathbf{R}^{\mathbf{N}}\) , let (for each \(n \in \mathbf{N}\) ) \(\mathfrak{m}_n\) be the maximal ideal consisting of the \(f\colon \mathbf{N} \to \mathbf{R}\) such that \(f(n) = 0\) . Show that \(\mathfrak{a} = \mathbf{R}^{(\mathbf{N})}\) is an ideal of \(\mathbf{A}\) contained in the union of the \(\mathfrak{m}_n\) , but in none of these maximal ideals.
\item
  In a ring A let p be a prime ideal and a an element such that \(p \subset Aa\) but \(a \notin p\) . Show that p = ap.
\item
  \begin{enumerate}
  \def\labelenumii{(\alph{enumii})}
  \tightlist
  \item
    Let A be a Noetherian ring and \(\mathfrak{a}\) an ideal of A. Show that there exists a finite number of prime ideals \(\mathfrak{p}_i\) ( \(1 \leqslant i \leqslant r\) ) such that \(\mathfrak{p}_1\mathfrak{p}_2 \ldots \mathfrak{p}_r \subset \mathfrak{a}\) . (Argue by reductio ad absurdum considering among the ideals containing \(\mathfrak{a}\) , distinct from A and containing no finite product of prime ideals, a maximal element;
  \end{enumerate}
\end{enumerate}

observe on the other hand that, if the ideal b is not prime, there exist two ideals c, b containing b, distinct from b and satisfying cd ⊂ b.) (Cf. Chapter IV.)

If A is an integral domain and \(a \neq 0\) , we may assume that the \(p_{i} \neq 0\) .

\begin{enumerate}
\def\labelenumi{(\alph{enumi})}
\setcounter{enumi}{1}
\tightlist
\item
  Give an example of a non-Noetherian ring A for which the result of (a) is false. (Take for example A to be the ring of continuous real-valued functions on {[}0, 1{]}.)
\end{enumerate}

¶ 6. (a) Let A be a ring and a, b two ideals of A such that b is finitely generated; show that, if the quotient rings A/a and A/b are Noetherian, so is A/ab (observe that b/ab is a finitely generated (A/a)-module).

\begin{enumerate}
\def\labelenumi{(\alph{enumi})}
\setcounter{enumi}{1}
\tightlist
\item
  Show that, if a ring A is such that every prime ideal of A is finitely generated, A is Noetherian. (Argue by reductio ad absurdum by showing that in the set of ideals of A which are not finitely generated there would be a maximal element c which is not prime by hypothesis; there would therefore be two ideals \(a \supset c\) , \(b \supset c\) distinct from c and satisfying \(ab \subset c\) ; then use (a).)
\end{enumerate}

\begin{enumerate}
\def\labelenumi{\arabic{enumi}.}
\setcounter{enumi}{6}
\tightlist
\item
  Let \(\mathbf{A} = \prod_{i=1}^{n} \mathbf{A}_i\) be the product of a finite family of rings, the \(\mathbf{A}_i\) being canonically identified with ideals of \(\mathbf{A}\) . Let \(\mathbf{B}\) be a subring of \(\mathbf{A}\) such that \(\operatorname{pr}_i \mathbf{B} = \mathbf{A}_i\) for \(1 \leqslant i \leqslant n\) .
\end{enumerate}

\begin{enumerate}
\def\labelenumi{(\alph{enumi})}
\item
  Show that, if the \(\mathbf{A}_i\) are Noetherian (resp. Artinian), B is Noetherian (resp. Artinian) (cf.~Algebra, Chapter VIII, § 2, Exercise 12).
\item
  Let \(n_{i}\) be the ideal of B which is the kernel of the restriction of \(pr_{i}\) to B. Show that every prime ideal p of B is contained in one of the \(n_{i}\) (use Proposition 1); deduce that, for this index i, \(pr_{i}\) \(p \neq A_{i}\) . Show that, if each of the \(A_{i}\) contains only a finite number \(k_{i}\) of maximal ideals, B contains at most \(\sum_{i=1}^{n} k_{i}\) maximal ideals.
\item
  Let \(\Re(A)\) , \(\Re(B)\) be the Jacobson radicals of \(A\) and \(B\) respectively. Show that \(\Re(B) = B \cap \Re(A)\) . If each of the \(A_i\) contains only a finite number of maximal ideals, show that, for every integer \(k > 1\) , \(\mathrm{pr}_i((\Re(B))^k) = (\Re(A_i))^k\) for all \(i\) and \((\Re(B))^k = (\Re(A))^k \cap B\) (write \(\Re(B)\) as a product of distinct maximal ideals and note that, if \(m\) is a maximal ideal of \(B\) such that \(\mathrm{pr}_i(m) = m_i\) is a maximal ideal of \(A_i\) , then \(\mathrm{pr}_i(m^k) = m_i^k\) ).
\end{enumerate}

\begin{enumerate}
\def\labelenumi{\arabic{enumi}.}
\setcounter{enumi}{7}
\tightlist
\item
  \begin{enumerate}
  \def\labelenumii{(\alph{enumii})}
  \tightlist
  \item
    Let \(\mathfrak{a}, \mathfrak{b}\) be two relatively prime ideals of \(\mathfrak{a}\) ring \(\mathbf{A}\) . Show that \(\mathfrak{a}: \mathfrak{b} = \mathfrak{a}\) , \(\mathfrak{b}: \mathfrak{a} = \mathfrak{b}\) . If \(\mathfrak{c}\) is an ideal of \(\mathbf{A}\) such that \(\mathfrak{ac} \subset \mathfrak{b}\) , then \(\mathfrak{c} \subset \mathfrak{b}\) .
  \end{enumerate}
\end{enumerate}

\begin{enumerate}
\def\labelenumi{(\alph{enumi})}
\setcounter{enumi}{1}
\item
  Let p, q be two prime ideals neither of which is contained in the other; then p: q = p and q: p = q. Give an example of two principal prime ideals p, q in the polynomial ring A = K{[}X, Y{]} (where K is a field), neither of which is contained in the other and which are not relatively prime.
\item
  Let \(a\) be an element which is not a divisor of 0 in A. Show that, if the principal ideal \(p = Aa\) is prime, the relation \(p = bc\) for two ideals \(b\) , \(c\) of A implies \(b = A\) or \(c = A\) .
\end{enumerate}

\begin{enumerate}
\def\labelenumi{\arabic{enumi}.}
\setcounter{enumi}{8}
\tightlist
\item
  \begin{enumerate}
  \def\labelenumii{(\alph{enumii})}
  \tightlist
  \item
    Let \((\mathfrak{a}_i)_{1 \leqslant i \leqslant n}\) be a finite family of ideals of a ring A which are relatively prime in pairs and such that no \(\mathfrak{a}_i\) is of the form \(\mathfrak{bc} = \mathfrak{b} \cap \mathfrak{c}\) , where \(\mathfrak{b}\) and \(\mathfrak{c}\) are two relatively prime ideals. If \((\mathfrak{a}_j')_{1 \leqslant j \leqslant m}\) is another family of ideals of A which are relatively prime in pairs, such that \(\mathfrak{a}_1\mathfrak{a}_2 \ldots \mathfrak{a}_n = \mathfrak{a}_1'\mathfrak{a}_2' \ldots \mathfrak{a}_m'\) , show that \(m \leqslant n\) ; if \(m = n\) , there exists a permutation \(\pi\) on \([1, n]\) such that \(a_i' = a_{\pi(i)}\) for \(1 \leqslant i \leqslant n\) (use Proposition 5 and Algebra, Chapter VIII, § 1, Exercise 1 (d)).
  \end{enumerate}
\end{enumerate}

\begin{enumerate}
\def\labelenumi{(\alph{enumi})}
\setcounter{enumi}{1}
\tightlist
\item
  Let A be a Noetherian ring. Show that every ideal \(\mathfrak{a}\) of A is equal to the product of a finite number of ideals which are relatively prime in pairs, none of which is the product of two relatively prime ideals. (Argue by reductio ad absurdum considering a maximal element of the set of ideals not having this property.)
\end{enumerate}

\begin{enumerate}
\def\labelenumi{\arabic{enumi}.}
\setcounter{enumi}{9}
\tightlist
\item
  Give an example of a ring A and an infinite family of distinct maximal ideals \((\mathfrak{m}_n)_{n\in \mathbf{N}}\) of A such that \(\bigcap_{n}\mathfrak{m}_{n} = 0\) but the canonical mapping
\end{enumerate}

\[
\mathrm{A} \rightarrow \prod_ {n} (\mathrm{A} / \mathrm{m} _ {n})
\]

is not surjective.

\begin{enumerate}
\def\labelenumi{\arabic{enumi}.}
\setcounter{enumi}{10}
\tightlist
\item
  Let A be a ring which is not necessarily commutative and let \(a, b\) be two two-sided ideals of A such that \(a + b = A\) . Show that \(a \cap b = ab + ba\) . Give an example where \(a \cap b \neq ab\) (consider the ring of lower triangular matrices of order 2 over a field).
\end{enumerate}

\exercisesection{2}

\begin{enumerate}
\def\labelenumi{\arabic{enumi}.}
\tightlist
\item
  A multiplicative subset \(S\) of a ring \(A\) is called saturated if the relation \(xy \in S\) implies \(x \in S\) and \(y \in S\) .
\end{enumerate}

\begin{enumerate}
\def\labelenumi{(\alph{enumi})}
\item
  For a subset \(S\) of \(A\) to be multiplicative and saturated, it is necessary and sufficient that \(A - S\) be a union of prime ideals of \(A\) .
\item
  Let S be a multiplicative subset of A and \(\tilde{S}\) the set of \(x \in A\) for which there exists \(y \in A\) such that \(xy \in S\) . Show that \(\tilde{S}\) is the smallest saturated multiplicative subset containing S, A --- \(\tilde{S}\) is the union of the prime ideals of A not meeting S and, for every A-module M, the canonical mapping from \(S^{-1}M\) to \(\tilde{S}^{-1}M\) is bijective.
\item
  Let S and T be two multiplicative subsets of A. Show that the two following properties are equivalent: (α) \(S \subset \tilde{T}\) ; (β) for every A-module M such that \(S^{-1}M = 0\) , \(T^{-1}M = 0\) . (To see that (β) implies (α), consider a quotient A-module A/As, where \(s \in S\) .)
\end{enumerate}

\begin{enumerate}
\def\labelenumi{\arabic{enumi}.}
\setcounter{enumi}{1}
\tightlist
\item
  Let A be a ring and S a multiplicative subset of A. Show that the set n of \(x \in A\) such that sx = 0 for some \(s \in S\) is an ideal of A. Set \(A_{1} = A/n\) and denote by \(S_{1}\) the canonical image of S in \(A_{1}\) . Show that no element of \(S_{1}\) is a divisor of 0 in \(\mathbf{A}_1\) and that the canonical homomorphism \(\mathbf{S}^{-1}\mathbf{A} \to \mathbf{S}_1^{-1}\mathbf{A}_1\) is bijective.
\end{enumerate}

Deduce that, for \(S^{-1}A\) to be a finitely generated A-module, it is necessary and sufficient that all the elements of \(S_{1}\) be invertible in \(A_{1}\) , in which case \(S^{-1}A\) is identified with \(A_{1}\) .

\begin{enumerate}
\def\labelenumi{\arabic{enumi}.}
\setcounter{enumi}{2}
\tightlist
\item
  Let \(S = Z^{*}\) (the complement of \(\{0\}\) in Z), p an integer \textgreater1 and M the Z-module the direct sum of the modules \(Z/p^{n}Z\) for \(n \in N\) . Show that \(S^{-1}M = 0\) although M is a faithful Z-module and that the canonical mapping
\end{enumerate}

\[
\mathrm{S} ^ {- 1} \operatorname{End} _ {\mathbf {z}} (\mathrm{M}) \rightarrow \operatorname{End} _ {\mathrm{s} ^ {- 1} \mathbf {z}} (\mathrm{S} ^ {- 1} \mathrm{M})
\]

is not injective.

\begin{enumerate}
\def\labelenumi{\arabic{enumi}.}
\setcounter{enumi}{3}
\item
  Let \(S = \mathbf{Z}^*\) and \(M\) be a free \(\mathbf{Z}\) -module with infinite basis. Show that the canonical mapping \(S^{-1}\operatorname{End}_{\mathbf{Z}}(M) \to \operatorname{End}_{S^{-1}\mathbf{Z}}(S^{-1}M)\) is not surjective. Deduce from this and Exercise 3 an example of a \(\mathbf{Z}\) -module \(M\) such that the canonical mapping \(S^{-1}\operatorname{End}_{\mathbf{Z}}(M) \to \operatorname{End}_{S^{-1}\mathbf{Z}}(S^{-1}M)\) is neither injective nor surjective.
\item
  Give an example of a sequence \((\mathbf{P}_n)\) of sub- \(\mathbf{Z}\) -modules of \(\mathbf{Z}\) such that, for \(\mathbf{S} = \mathbf{Z}^*\) , the submodule \(\mathrm{S}^{-1}\left(\bigcap_{n} \mathrm{P}_{n}\right)\) is distinct from \(\bigcap_{n} (\mathrm{S}^{-1}\mathrm{P}_{n})\) .
\item
  Give an example of a \(\mathbf{Z}\) -module \(\mathbf{M}\) such that, for \(\mathbf{S} = \mathbf{Z}^*\) , \(\operatorname{Ann}(\mathbf{M}) = 0\) , but \(\operatorname{Ann}(\mathbf{S}^{-1}\mathbf{M}) = \mathbf{Q}\) (cf.~Exercise 3).
\item
  Let S be a multiplicative subset of a ring A; for every family \((M_{t})_{t\in I}\) of A-modules, define a canonical homomorphism \(S^{-1}\prod_{i\in I}M_{i}\to\prod_{i\in I}S^{-1}M_{i}\) of \(S^{-1}A\) -modules. Give an example where this homomorphism is neither injective nor surjective (cf.~Exercise 3).
\item
  Let A be a ring, \((\mathrm{S}_{\lambda})_{\lambda \in L}\) a family of multiplicative subsets of A and \(M = \bigoplus_{\lambda \in L} S_{\lambda}^{-1}A\) . Show that the two following properties are equivalent: ( \(\alpha\) ) M is a faithfully flat A-module; ( \(\beta\) ) for every maximal ideal m of A, there exists \(\lambda \in L\) such that \(m \cap S_{\lambda} = \varnothing\) . In particular, if S is a multiplicative subset of A, the A-module \(S^{-1}A\) is faithfully flat only if the elements of S are invertible, in which case \(S^{-1}A\) is identified with A.
\item
  Let A be a ring, S a multiplicative subset of A and q a prime ideal of A. Show that, if T is the image of A - q in \(S^{-1}A\) , the ring \(T^{-1}(S^{-1}A)\) is isomorphic to \(S'^{-1}A\) , where \(S'\) is the complement of the union of the prime ideals of A contained in q and not meeting S.
\item
  \begin{enumerate}
  \def\labelenumii{(\alph{enumii})}
  \tightlist
  \item
    Let K be a commutative field, A the polynomial ring K{[}X, Y{]}, \(p = (X)\) , \(q = (Y)\) and S the complement of \(p \cup q\) in A. Show that the saturation of \(p + q\) with respect to S is distinct from the sum of the saturations of p and q.
  \end{enumerate}
\end{enumerate}

\begin{enumerate}
\def\labelenumi{(\alph{enumi})}
\setcounter{enumi}{1}
\tightlist
\item
  Let K be a commutative field, A the polynomial ring K{[}X, Y, Z{]}, \(\mathfrak{p} = (\mathbf{X}) + (\mathbf{Z})\) , \(\mathfrak{q} = (\mathbf{Y}) + (\mathbf{Z})\) and S the complement of \(p \cup q\) in A. Show that the saturation of pq with respect to S is distinct from the product of the saturations of p and q.
\end{enumerate}

\begin{enumerate}
\def\labelenumi{\arabic{enumi}.}
\setcounter{enumi}{10}
\item
  Let \(p\) be a prime number; which ideals of \(\mathbf{Z}\) are saturated with respect to the ideal \((p)\) ?
\item
  In a ring A let \(\mathfrak{r}(\mathfrak{a})\) denote the radical of an ideal \(\mathfrak{a}\) .
\end{enumerate}

\begin{enumerate}
\def\labelenumi{(\alph{enumi})}
\tightlist
\item
  Show that, for three ideals a, b, c of A,
\end{enumerate}

\[
\mathfrak {r} (\mathfrak {a} + \mathfrak {b c}) = \mathfrak {r} (\mathfrak {a} + (\mathfrak {b} \cap \mathfrak {c})) = \mathfrak {r} (\mathfrak {a} + \mathfrak {b}) \cap \mathfrak {r} (\mathfrak {a} + \mathfrak {c}).
\]

\begin{enumerate}
\def\labelenumi{(\alph{enumi})}
\setcounter{enumi}{1}
\tightlist
\item
  Give an example of an infinite sequence of ideals \((\mathfrak{a}_n)\) of A such that \(\mathfrak{r}\left(\bigcap_{n} \mathfrak{a}_n\right) \neq \bigcap_{n} \mathfrak{r}(\mathfrak{a}_n)\) .
\end{enumerate}

*(c) Give an example of a principal ideal \(\mathfrak{a}\) such that \(\mathfrak{r}(\mathfrak{a})\) is not finitely generated and there exists no integer \(n > 0\) such that \((\mathfrak{r}(\mathfrak{a}))^n \subset \mathfrak{a}\) . (Consider the ring of a non-discrete valuation of height 1.)*

13. Let A be a ring.\par

\begin{enumerate}
\def\labelenumi{(\alph{enumi})}
\item
  Give another proof of Exercise 6 (b) of Algebra, Chapter VIII, § 6, by considering the image of the polynomial \(f\) in the rings \((\mathbf{A} / \mathfrak{p})[\mathbf{X}]\) , where \(\mathfrak{p}\) runs through the set of prime ideals of \(\mathbf{A}\) .
\item
  Let \(N\) be a square matrix of order \(r\) over \(A\) ; for \(N\) to be nilpotent, it is necessary and sufficient that each of the coefficients (other than the dominant coefficient) of its characteristic polynomial be so. (To see that the condition is necessary, use the same method as in (a); to see that it is sufficient, use the Cayley-Hamilton Theorem.)
\item
  Show that, for every ordered pair of positive integers \(k, r\) , there exists a least positive integer \(m(k, r)\) independent of the ring \(\mathbf{A}\) such that the relation \(N^k = 0\) implies \((\operatorname{Tr}(N))^{m(k,r)} = 0\) . (Consider the ring \(\Lambda = \mathbf{Z}[T_{ij}]\) , where the \(T_{ij}\) are \(r^2\) indeterminates, the matrix \(X = (T_{ij})\) over \(\Lambda\) and the element \(t = \sum_{i=1}^{n} \mathrm{T}_{ii} = \mathrm{Tr}(X)\) of \(\Lambda\) ; if \(\mathfrak{a}_k\) is the ideal of \(\Lambda\) generated by the elements of \(X^k\) , show, with the aid of (b), that there exists an integer \(m\) such that \(t^m \in \mathfrak{a}_k\) .)
\item
  Show that \(m(2, 2) = 4\) .
\end{enumerate}

¶ 14. (a) Let A be a left Noetherian ring (not necessarily commutative) and a a left ideal of A. Suppose that there exists a left ideal b ≠ 0 such that a ∩ b = 0.

Show that every element \(a \in \mathfrak{a}\) is a right divisor of 0. (Let \(b \neq 0\) be an element of b and let \(a_{n} = Ab + Aba + \cdots + Aba^{n}\) ; consider the smallest integer n such that \(a_{n} = a_{n+1}\) .

\begin{enumerate}
\def\labelenumi{(\alph{enumi})}
\setcounter{enumi}{1}
\tightlist
\item
  Let E be an infinite-dimensional vector space over a field K and A the ring \(\operatorname{End}_{\mathbb{K}}(\operatorname{E})\) . Give an example of two elements u, v of A such that neither u nor v is a right divisor of 0 and \(Au \cap Av = 0\) .
\end{enumerate}

¶ 15. Let X be a completely regular topological space and βX its Stone-Čech compactification (General Topology, Chapter IX, § 1, Exercise 7). For all \(x \in \mathbf{X}\) , let \(\mathfrak{J}(x)\) denote the maximal ideal of the ring \(\mathcal{C}(\mathbf{X}; \mathbf{R})\) consisting of those \(f \in \mathcal{C}(\mathbf{X}; \mathbf{R})\) such that \(x\) belongs to the closure of the sub set \(\bar{f}^1(0)\) of X (General Topology, Chapter X, § 4, Exercise 15); let \(\mathfrak{U}(x)\) denote the ideal consisting of those \(f \in \mathcal{C}(\mathbf{X}; \mathbf{R})\) such that \(\bar{f}^1(0)\) is the trace on X of a neighbourhood of \(x\) in X; \(\mathfrak{U}(x)\) is equal to its radical.

An ideal \(\mathfrak{a}\) of \(\mathfrak{C}(\mathbf{X};\mathbf{R})\) is called isolated (resp. absolutely isolated) if, for all \(f\geqslant 0\) on \(\mathfrak{a}\) , every \(g\) such that \(0\leqslant g\leqslant f\) (resp. \(|g|\leqslant f\) ) belongs to \(\mathfrak{a}\) (cf.~Algebra, Chapter IV, §1, Exercise 4); then \(\mathcal{C}(\mathbf{X};\mathbf{R}) / \mathfrak{a}\) possesses an order structure (resp. a lattice-order structure) compatible with its ring structure and for which the elements \(\geqslant 0\) are the classes mod. \(\mathfrak{a}\) of the elements \(\geqslant 0\) to \(\mathcal{C}(\mathbf{X};\mathbf{R})\) (loc. cit.).

\begin{enumerate}
\def\labelenumi{(\alph{enumi})}
\item
  Show that an ideal \(\mathfrak{a}\) of \(\mathcal{C}(\mathbf{X};\mathbf{R})\) which is equal to its radical is absolutely isolated (if \(f\in \mathfrak{a}\) and \(|g|\leqslant |f|\) , consider the function equal to 0 when \(f(x) = 0\) and to \((g(x))^2 /f(x)\) when \(f(x)\neq 0\) ). In this case for \(\mathcal{C}(\mathbf{X};\mathbf{R}) / \mathfrak{a}\) to be totally ordered, it is necessary and sufficient that \(\mathfrak{a}\) be prime; the set of prime ideals of \(\mathcal{C}(\mathbf{X};\mathbf{R}) / \mathfrak{a}\) is then totally ordered by inclusion.
\item
  For an isolated ideal \(\mathfrak{a}\) of \(\mathcal{C}(\dot{\mathbf{X}};\mathbf{R})\) to be such that \(\mathcal{C}(\mathbf{X};\mathbf{R}) / \mathfrak{a}\) is totally ordered, it is necessary that \(\mathfrak{a}\) be contained in a single maximal ideal \(\mathfrak{J}(x)\) and then \(\mathfrak{a}\) necessarily contains \(\mathfrak{U}(x)\) ; it is sufficient that \(\mathfrak{a}\) contain a prime ideal and then \(\mathfrak{a}\) is absolutely isolated (note that the relation \(fg = 0\) then implies \(f\in \mathfrak{a}\) or \(g\in \mathfrak{a}\) ).
\item
  Take \(\mathbf{X} = \mathbf{R}\) , \(x = 0\) ; show that the principal ideal generated by the function \(t \mapsto |t|\) is contained in \(\mathfrak{J}(0)\) and contains \(\mathfrak{U}(0)\) but is not isolated; the principal ideal generated by the identity mapping \(t \mapsto t\) is contained in \(\mathfrak{J}(0)\) , contains \(\mathfrak{U}(0)\) and is isolated but not absolutely isolated.
\item
  Show that, if \(\mathfrak{p}\) and \(\mathfrak{q}\) are two prime ideals in \(\mathcal{C}(\mathbf{X};\mathbf{R})\) , then \(\mathfrak{pq} = \mathfrak{p} \cap \mathfrak{q}\) (and in particular \(\mathfrak{p}^2 = \mathfrak{p}\) ) and \(\mathfrak{p} + \mathfrak{q}\) is prime or equal to \(\mathcal{C}(\mathbf{X};\mathbf{R})\) . (If \(\mathfrak{p} + \mathfrak{q} \neq \mathcal{C}(\mathbf{X};\mathbf{R})\) , observe that \(\mathfrak{p} + \mathfrak{q}\) is absolutely isolated and equal to its radical and use (a).)
\item
  For \(\mathfrak{U}(x)\) to be prime, it is necessary and sufficient that, for every function \(f\in\mathcal{C}(\mathbf{X};\mathbf{R})\) , the closure in \(\beta\mathbf{X}\) of one of the two sets \(\bar{f}^{-1}\left([0,+\infty[),\bar{f}^{-1}([- \infty,0])\right)\) is a neighbourhood of \(x\) . Deduce an example where \(\mathfrak{U}(x)\) is prime and not maximal (consider the topological space associated with a non-trivial ultrafilter; cf.~General Topology, Chapter I, § 6, no. 5).
\item
  Suppose that \(x \in \mathbf{X}\) . For \(\mathfrak{U}(x) = \mathfrak{J}(x)\) , it is necessary and sufficient that every countable intersection of neighbourhoods of \(x\) in \(\mathbf{X}\) be a neighbourhood of \(x\) in \(\mathbf{X}\) . (Cf. General Topology, Chapter I, § 7, Exercise 7.)
\item
  Show that \(\mathcal{C}(\mathbf{X};\mathbf{R})\) is canonically identified with the total ring of fractions of \(\mathcal{C}^{\infty}(\mathbf{X};\mathbf{R})\) .
\end{enumerate}

\begin{enumerate}
\def\labelenumi{\arabic{enumi}.}
\setcounter{enumi}{15}
\tightlist
\item
  Let A be a topological ring (not necessarily commutative). The topology on A is called (left) linear if the open left ideals of A form a fundamental system of neighbourhoods of 0. If this is so, the set F of open left ideals of A satisfies the following conditions:
\end{enumerate}

\begin{enumerate}
\def\labelenumi{(\arabic{enumi})}
\item
  Every left ideal of A containing an ideal \(m \in \mathcal{F}\) belongs to \(\mathcal{F}\) .
\item
  Every finite intersection of ideals of \(\mathcal{F}\) belongs to \(\mathcal{F}\) .
\item
  If \(\mathfrak{m} \in \mathcal{F}\) and \(a \in \mathbf{A}\) , the left ideal \(\mathfrak{n}\) consists of those \(x \in \mathbf{A}\) such that \(xa \in \mathfrak{m}\) belongs to \(\mathcal{F}\) .
\end{enumerate}

Conversely, if a set F of left ideals of a ring A satisfies these three conditions, there exists a unique topology compatible with the ring structure on A and for which F is the set of open left ideals of A. Such a set F of left ideals is called (left) topologizing.

\begin{enumerate}
\def\labelenumi{\arabic{enumi}.}
\setcounter{enumi}{16}
\tightlist
\item
  (*) Let A be a ring (not necessarily commutative) and \(\mathcal{F}\) a topologizing set (Exercise 16) of left ideals of A.
\end{enumerate}

\begin{enumerate}
\def\labelenumi{(\alph{enumi})}
\item
  A (left) A-module M is called F-negligible if the annihilator of every element of M belongs to F. If M is F-negligible, so is every submodule and every quotient module of M. If M and N are two F negligible A-modules, \(M \oplus N\) is F-negligible. For \(A_{s}/m\) to be F-negligible, it is necessary and sufficient that \(m \in F\) . If \((0) \in F\) , every A-module is F-negligible; if consists only of the single left ideal \(A_{s}\) , every F-negligible A-module is reduced to 0.
\item
  For every A-module M there exists a greatest F-negligible submodule of M which is denoted by FM. For every A-module homomorphism \(u: M \to N\) , \(u(\mathcal{F}M) \subset \mathcal{F}N\) ; let Fu be the mapping from FM to FN whose graph coincides with that of \(u|FM\) . If
\end{enumerate}

\[
0 \longrightarrow \mathrm{M} \stackrel {u} {\longrightarrow} \mathrm{N} \stackrel {v} {\longrightarrow} \mathrm{P}
\]

is an exact sequence of A-modules, the sequence

\[
0 \longrightarrow \mathcal {F} M \stackrel {\mathcal {F} _ {u}} {\longrightarrow} \mathcal {F} N \stackrel {\mathcal {F} _ {v}} {\longrightarrow} \mathcal {F} P
\]

is exact. Show that \(FA_{s}\) is a two-sided ideal of A.

\begin{enumerate}
\def\labelenumi{(\alph{enumi})}
\setcounter{enumi}{2}
\tightlist
\item
  Let \(\mathfrak{m}\) , \(\mathfrak{n}\) be two elements of \(\mathcal{F}\) such that \(\mathfrak{m} \subset \mathfrak{n}\) ; for every left A-module \(\mathbf{M}\) the canonical injection \(j \colon \mathfrak{m} \to \mathfrak{n}\) defines a commutative group homomorphism
\end{enumerate}

\[
u _ {m, n} = \operatorname{Hom} _ {A} (j, 1 _ {M}): \operatorname{Hom} _ {A} (n, M) \rightarrow \operatorname{Hom} _ {A} (m, M).
\]

If \(\mathcal{F}\) is ordered by the relation \(\supset\) , the \(u_{m,n}\) define a direct system of commutative groups, whose direct limit will be denoted by \(\mathrm{M}_{(\mathcal{F})}\) . The mappings \(u_{m,A_s}\) form a direct system, whose direct limit is a mapping (called canonical) from \(\mathrm{M} = \mathrm{Hom}_{\mathrm{A}}(\mathrm{A}_s,\mathrm{M})\) to \(\mathrm{M}_{(\mathcal{F})}\) . Similarly, if \(v: \mathrm{M} \to \mathrm{N}\) is a left A-module homomorphism, the \(v_m = \mathrm{Hom}_{\mathrm{A}}(1,v): \mathrm{Hom}_{\mathrm{A}}(\mathfrak{m},\mathrm{M}) \to \mathrm{Hom}_{\mathrm{A}}(\mathfrak{m},\mathrm{N})\) form a direct system, whose direct limit is a commutative group homomorphism

\[
v _ {(\mathcal {F})} \colon \mathrm{M} _ {(\mathcal {F})} \rightarrow \mathrm{N} _ {(\mathcal {F})}.
\]

If \(0 \to \mathbf{M} \xrightarrow{v} \mathbf{N} \xrightarrow{w} \mathbf{P}\) is an exact sequence of A-modules, the sequence

\[
0 \longrightarrow \mathrm{M} _ {(\mathcal {F})} \stackrel {v (\mathcal {F})} {\longrightarrow} \mathrm{N} _ {(\mathcal {F})} \stackrel {w (\mathcal {F})} {\longrightarrow} \mathrm{P} _ {(\mathcal {F})}
\]

is exact.

\begin{enumerate}
\def\labelenumi{\arabic{enumi}.}
\setcounter{enumi}{17}
\tightlist
\item
  \begin{enumerate}
  \def\labelenumii{(\alph{enumii})}
  \tightlist
  \item
    Let A be a ring and F and G two topologizing sets (Exercise 16) of left ideals of A. Let F.G denote the set of left ideals l of A for which there exists a left ideal n ∈ G such that n/l is F-negligible. Show that F.G is topologizing; for an A-module M to be F.G-negligible, it is necessary and sufficient that there exist an F-negligible submodule M' of M such that M/M' is G-negligible. Show that the composition law (F,G) → F.G on the set of topologizing sets of left ideals of A is associative.
  \end{enumerate}
\end{enumerate}

\(\mathcal{F}\) is called idempotent if \(\mathcal{F}.\mathcal{F} = \mathcal{F}\) .

\begin{enumerate}
\def\labelenumi{(\alph{enumi})}
\setcounter{enumi}{1}
\item
  Show that \(\mathcal{F}.\mathcal{G}\) contains all the ideals of the form \(\mathfrak{m}.\mathfrak{n}\) where \(\mathfrak{m}\in\mathcal{F}\) and \(\mathfrak{n}\in\mathcal{G}\) .
\item
  Let K be a commutative field and B the ring of polynomials with coefficients in K in an infinite system of indeterminates \((\mathbf{X}_{i})\) ( \(i \geqslant 0\) ). Let b be the ideal of B generated by the elements \(X_{i}X_{j}\) ( \(i \neq j\) ); let A be the ring B/b, \(\xi_{i}\) the class mod. b of \(X_{i}\) and m the ideal of A generated by the \(\xi_{i}\) ; take F to be the set of ideals containing a power of m. Show that F is topologizing and contains the product of any two ideals of F but is not idempotent.
\item
  Suppose that A is commutative and F is a topologizing set of ideals of A such that every ideal \(m \in F\) contains a finitely generated ideal \(n \in F\) . Show that, if the product of two ideals of F belongs to F, F is idempotent.
\end{enumerate}

¶ 19. Let F be a set of left ideals of a ring A; suppose that F is topologizing and idempotent (Exercises 16 and 18).

\begin{enumerate}
\def\labelenumi{(\alph{enumi})}
\tightlist
\item
  Show that, if \(\mathfrak{m} \in \mathcal{F}\) , \(\mathfrak{n} \in \mathcal{F}\) and \(u \in \operatorname{Hom}_{\mathbf{A}}(\mathfrak{n}, \mathbf{A}_s)\) , then \(\bar{u}^1(\mathfrak{m}) \in \mathcal{F}\) (consider the exact sequence \(0 \to \mathfrak{n} / \bar{u}^1(\mathfrak{m}) \to \mathbf{A}_s / \bar{u}^1(\mathfrak{m}) \to \mathbf{A}_s / \mathfrak{n} \to 0\) ). For every A-module M and all \(v \in \operatorname{Hom}_{\mathbf{A}}(\mathfrak{m}, \mathbf{M})\) , let \(u.v\) denote the canonical image in \(\mathbf{M}_{(\mathcal{F})}\) (Exercise 17 (c)) of the composite homomorphism \(\bar{u}^1(\mathfrak{m}) \xrightarrow{u} \mathfrak{m} \xrightarrow{v} \mathbf{M}\) . Show that the mapping \((u,v) \mapsto u.v\) is Z-bilinear; there corresponds to it a Z-linear mapping
\end{enumerate}

\[
\psi_ {m, n} \colon \operatorname{Hom} _ {A} (n, A _ {s}) \otimes_ {\mathbf {Z}} \operatorname{Hom} _ {A} (m, M) \rightarrow M _ {(\mathcal {F})}.
\]

These mappings form a direct system and their direct limit therefore corresponds canonically to a Z-bilinear mapping

\[
\psi \colon A _ {(\mathcal {F})} \times M _ {(\mathcal {F})} \rightarrow M _ {(\mathcal {F})}.
\]

Show that, if \(M = A_{s}\) , \(\psi\) is an internal law of composition which makes \(\mathbf{A}_{(\mathcal{F})}\) into a ring and the canonical mapping \(A \to A_{(\mathcal{F})}\) is a ring homomorphism for this structure. For any left A-module M, \(\psi\) is an external law defining on \(\mathbf{M}_{(\mathcal{F})}\) a left \(\mathbf{A}_{(\mathcal{F})}\) -module structure; with this structure (and the canonical mapping \(A \to A_{(\mathcal{F})}\) ) the canonical mapping \(M \to M_{(\mathcal{F})}\) is an A-module homomorphism.

\begin{enumerate}
\def\labelenumi{(\alph{enumi})}
\setcounter{enumi}{1}
\tightlist
\item
  If \(i: \mathbf{M} \to \mathbf{M}_{(\mathcal{F})}\) is the canonical homomorphism, show that
\end{enumerate}

\[
\operatorname{Ker} (i) = \mathcal {F M}
\]

and that the A-module Coker \((i) = \mathbf{M}_{(\mathcal{F})} / i(\mathbf{M})\) is \(\mathcal{F}\) -negligible.

\begin{enumerate}
\def\labelenumi{(\alph{enumi})}
\setcounter{enumi}{2}
\tightlist
\item
  For every left A-module M let \(M_{F}\) denote the left A-module \((\mathrm{M}/\mathcal{F}\mathrm{M})_{(\mathcal{F})}\) and \(j_{M}\) the composite canonical mapping
\end{enumerate}

\[
\mathbf {M} \rightarrow \mathbf {M} / \mathcal {F} \mathbf {M} \rightarrow \mathbf {M} _ {\mathcal {F}};
\]

then \(\operatorname{Ker}(j_{\mathbf{M}}) = \mathcal{F}\mathbf{M}\) and the A-module \(\operatorname{Coker}(j_{\mathbf{M}}) = \mathrm{M}_{\mathcal{F}} / j_{\mathbf{M}}(\mathbf{M})\) is \(\mathcal{F}\) -negligible.

Let \(u: \mathrm{P} \to \mathrm{M}\) , \(h: \mathrm{P} \to \mathrm{Q}\) be A-module homomorphisms such that \(\operatorname{Ker}(h)\) and \(\operatorname{Coker}(h)\) are \(\mathcal{F}\) -negligible; show that there exists a unique A-homomorphism \(v: \mathrm{Q} \to \mathrm{M}_{\mathcal{F}}\) which makes the following diagram commutative

\[
\begin{array}{c} \mathrm{P} \xrightarrow {u} \mathrm{M} \\ h \Big \downarrow \qquad \qquad \qquad \qquad \qquad \qquad \qquad \Big \downarrow j _ {\mathrm{M}} \\ \mathrm{Q} \xrightarrow [ v ]{} \mathrm{M} _ {\mathcal {F}} \end{array}
\]

(First reduce it to the case where h is injective; then, identifying P with a submodule of Q, note that, for \(x \in Q\) , there exists \(m \in F\) such that \(mx \in P\) and map x to the canonical image in \(M_{F}\) of the composite homomorphism

\[
\mathfrak {m} \xrightarrow {x} \mathfrak {m} x \to \mathrm{P} \xrightarrow {u} \mathrm{M} \to \mathrm{M} / \mathcal {F} \mathrm{M.})
\]

\begin{enumerate}
\def\labelenumi{(\alph{enumi})}
\setcounter{enumi}{3}
\tightlist
\item
  For every A-module homomorphism \(u: M \rightarrow N\) show that there exists a unique homomorphism \(u_{F}: M_{F} \rightarrow N_{F}\) which makes the following diagram commutative
\end{enumerate}

\[
\begin{array}{c c c} \mathrm{M} & \stackrel {{u}} {{\longrightarrow}} \mathrm{N} \\ j _ {\mathrm{M}} \Big \downarrow & & \Big \downarrow j _ {\mathrm{N}} \\ \mathrm{M} _ {\mathcal {F}} & \stackrel {{u _ {\mathcal {F}}}} {{\longrightarrow}} \mathrm{N} _ {\mathcal {F}} \end{array}
\]

(use (a), (b) and (c)). Show that, if

\[
0 \longrightarrow \mathrm{M} \stackrel {u} {\longrightarrow} \mathrm{N} \stackrel {v} {\longrightarrow} \mathrm{P}
\]

is an exact sequence of A-modules, the sequence

\[
0 \longrightarrow \mathrm{M} _ {\mathcal {F}} \stackrel {u \mathcal {F}} {\longrightarrow} \mathrm{N} _ {\mathcal {F}} \stackrel {v \mathcal {F}} {\longrightarrow} \mathrm{P} _ {\mathcal {F}}
\]

is exact (note that the mapping \(\mathbf{M} / \mathcal{F}\mathbf{M}\to \mathrm{N} / \mathcal{F}\mathrm{N}\) derived from \(u\) by taking quotients is injective).

Show that the mappings \(j_{\mathbf{M}_{\mathcal{F}}}\) and \((j_{\mathbf{M}})_{\mathcal{F}}\) coincide and are isomorphisms (and therefore \(\mathcal{FM}_{\mathcal{F}} = 0\) ). If \(\mathbf{N}'\) is a sub-A-module of \(\mathbf{M}_{\mathcal{F}}\) , \((\overline{j_{\mathbf{M}}^{-1}(\mathbf{N}')})_{\mathcal{F}}\) is canonically identified with \(\mathbf{N}_{\mathcal{F}}'\) (observe that, if we write \(\mathbf{N} = j_{\mathbf{M}}^{-1}(\mathbf{N}')\) , \(\mathbf{N}' / j_{\mathbf{M}}(\mathbf{N})\) is \(\mathcal{F}\) -negligible).

\begin{enumerate}
\def\labelenumi{(\alph{enumi})}
\setcounter{enumi}{4}
\tightlist
\item
  Let M be a left A-module and \(\phi: \mathbf{A} \times \mathbf{M} \to \mathbf{M}\) the Z-bilinear mapping \((a, m) \mapsto am\) . Show that there exists a unique Z-bilinear mapping \(\phi_{\mathcal{F}}: \mathbf{A}_{\mathcal{F}} \times \mathbf{M}_{\mathcal{F}} \to \mathbf{M}_{\mathcal{F}}\) which makes the following diagram commutative
\end{enumerate}

\[
\begin{array}{c} \mathbf {A} \times \mathbf {M} \xrightarrow {\phi} \mathbf {M} \\ j _ {\mathbf {A}} \times j _ {\mathbf {M}} \Big \downarrow \qquad \qquad \qquad \qquad \Big \downarrow j _ {\mathbf {M}} \\ \mathbf {A} _ {\mathcal {F}} \times \mathbf {M} _ {\mathcal {F}} \xrightarrow [ \phi_ {\mathcal {F}} ]{} \mathbf {M} _ {\mathcal {F}} \end{array}
\]

We also define on \(A_{F}\) a ring structure and on \(M_{F}\) a left \(A_{F}\) -module structure. If \(u: M \to N\) is an A-module homomorphism, \(u_{F}: M_{F} \to N_{F}\) is an \(A_{F}\) -module homomorphism. In order that, for every A-module M, the mapping \(u \mapsto u_{F}\) of \(\operatorname{Hom}_{\mathbf{A}}(\mathbf{M}, \mathbf{N})\) to \(\operatorname{Hom}_{\mathbf{A}_{F}}(\mathbf{M}_{F}, \mathbf{N}_{F})\) be bijective, it is necessary and sufficient that \(j_{N}\) be bijective.

\begin{enumerate}
\def\labelenumi{(\alph{enumi})}
\setcounter{enumi}{5}
\item
  Take A to be the polynomial ring K{[}X, Y{]} over a commutative field K and F to be the (topologizing and idempotent) set of ideals of A containing a power of the ideal m = (X) + (Y). If M = A/AX and u: A → M is the canonical mapping, show that \(u_{F}: A_{F} \to M_{F}\) is not surjective. (Show that \(A_{F}\) is identified with A and \(M_{F}\) with K{[} \(\overline{Y}\) , 1/ \(\overline{Y}\) {]}, where \(\overline{Y}\) is the class of Y in M, so that X \(\overline{Y}\) = 0 and X(1/ \(\overline{Y}\) ) = 0.)
\item
  If we set \(\mathcal{F}^1\mathrm{M} = \mathrm{M}_{\mathcal{F}} / j_{\mathrm{M}}(\mathrm{M})\) , show that, for every exact sequence \(0\to \mathrm{M}\to \mathrm{N}\to \mathrm{P}\to 0\) of A-modules, there is an exact sequence
\end{enumerate}

\[
0 \rightarrow \mathcal {F} M \rightarrow \mathcal {F} N \rightarrow \mathcal {F} P \rightarrow \mathcal {F} ^ {1} M \rightarrow \mathcal {F} ^ {1} N \rightarrow \mathcal {F} ^ {1} P
\]

(use the snake diagram).

\begin{enumerate}
\def\labelenumi{(\alph{enumi})}
\setcounter{enumi}{7}
\tightlist
\item
  Let \(\mathcal{F}'\) be the set of left ideals \(l'\) of \(\mathbf{A}_{\mathcal{F}}\) such that \((\mathbf{A}_{\mathcal{F}})_s / l'\) is an \(\mathcal{F}\) -negligible A-module: show that \(\mathcal{F}'\) is the set of left ideals of \(\mathbf{A}_{\mathcal{F}}\) containing the \(j_{\mathbf{A}}(\mathfrak{m})\) , where \(\mathfrak{m} \in \mathcal{F}\) (observe that \(\mathfrak{m}_{\mathcal{F}} = \mathbf{A}_{\mathcal{F}}\) for all \(\mathfrak{m} \in \mathcal{F}\) , using the exact sequence of (d); consequently \(A_{\mathcal{F}}/j_{A}(m)\) is F-negligible). For every \(A_{F}\) -module \(M'\) , \(FM' = F'M'\) ( \(M'\) being considered as an A-module by means of \(j_{A}\) ), \(F'\) is topologizing and idempotent and \(M'_{F'}\) is canonically identified with \(M'_{F}\) .
\end{enumerate}

\begin{enumerate}
\def\labelenumi{\arabic{enumi}.}
\setcounter{enumi}{19}
\tightlist
\item
  Let \(\mathcal{F}\) be a topologizing and idempotent set of left ideals of a ring A.
\end{enumerate}

\begin{enumerate}
\def\labelenumi{(\alph{enumi})}
\item
  Suppose that every ideal of \(\mathcal{F}\) contains a finitely generated ideal belonging to \(\mathcal{F}\) . Let M be an A-module and \((\mathbf{N}_{\iota})_{\iota \in I}\) a right directed family of submodules of M, whose union is M. Show that \(\mathbf{M}_{\mathcal{F}}\) is the union of its sub- \(\mathbf{A}_{\mathcal{F}}\) -modules \((\mathbf{N}_{\iota})_{\mathcal{F}}\) . In particular, for every family \((\mathbf{M}_{\lambda})_{\lambda \in L}\) of A-modules, \(\left( \bigoplus_{\lambda \in L} \mathbf{M}_{\lambda} \right)_{\mathcal{F}}\) is canonically isomorphic to \(\bigoplus_{\lambda \in L} (\mathbf{M}_{\lambda})_{\mathcal{F}}\) .
\item
  Suppose that condition (a) holds and also that, for every surjective A-module homomorphism \(u \colon \mathbf{M} \to \mathbf{N}\) , \(u_{\mathcal{F}} \colon \mathbf{M}_{\mathcal{F}} \to \mathbf{N}_{\mathcal{F}}\) is surjective. Then, the canonical mapping \(\mathbf{A}_{\mathcal{F}} \otimes_{\mathbf{A}} \mathbf{M} \to \mathbf{M}_{\mathcal{F}}\) defined by the external law on \(\mathbf{M}_{\mathcal{F}}\) as an \(\mathbf{A}_{\mathcal{F}}\) -module is an \(\mathbf{A}_{\mathcal{F}}\) -module isomorphism (reduce it to the case where \(\mathbf{M}\) is free); the ring \(\mathbf{A}_{\mathcal{F}}\) is a flat left A-module and the left \(\mathbf{A}_{\mathcal{F}}\) -modules coincide with the left A-modules \(\mathbf{M}\) such that the mapping \(j_{\mathbf{M}} \colon \mathbf{M} \to \mathbf{M}_{\mathcal{F}}\) is bijective.
\item
  Suppose that A is an infinite product \(\biguplus_{i\in I}A_{i}\) of rings; let \(e_{i}\) denote the unit element of \(A_{i}\) and take F to be the set of left ideals of A containing the two-sided ideal \(a=\bigoplus_{i\in I}A_{i}\) (the \(A_{i}\) being canonically identified with ideals of A). Show that, for every left A-module M, \(M_{F}\) is identified with the product \(\prod_{i\in I}(e_{i}M)\) and in particular \(A_{F}=A\) . Deduce that, for every surjective A-homomorphism \(u: M\to N\) , \(u_{F}: M_{F}\to N_{F}\) is surjective, but give examples of A-modules M such that \(j_{M}\) is not bijective.
\item
  Suppose that condition (a) holds and also that A is commutative. Show that, if M is F-negligible, \(\mathbf{M}_{(\mathcal{F})}=0\) .
\end{enumerate}

\begin{enumerate}
\def\labelenumi{\arabic{enumi}.}
\setcounter{enumi}{20}
\tightlist
\item
  Let A be a commutative ring and \(\mathcal{F}\) a topologizing and idempotent set of ideals of A.
\end{enumerate}

\begin{enumerate}
\def\labelenumi{(\alph{enumi})}
\tightlist
\item
  Then the ring \(\mathbf{A}_{\mathcal{F}}\) is commutative (note that, if \(\mathfrak{m} \in \mathcal{F}\) and
\end{enumerate}

\[
u \in \operatorname{Hom} _ {\mathbf {A}} (\mathfrak {m}, \mathbf {A}),
\]

then \(u(\mathfrak{m}^2) \subset \mathfrak{m}\) ; if \(v\) is another element of \(\operatorname{Hom}_{\mathbf{A}}(\mathfrak{m}, \mathbf{A})\) , show that

\[
v (u (x y)) = u (v (x y))
\]

for \(x \in \mathfrak{m}, y \in \mathfrak{m}\) ).

\begin{enumerate}
\def\labelenumi{(\alph{enumi})}
\setcounter{enumi}{1}
\item
  Let \(\mathcal{F}'\) be the family of ideals \(l'\) of \(\mathbf{A}_{\mathcal{F}}\) such that \(\mathbf{A}_{\mathcal{F}} / l'\) is \(\mathcal{F}\) -negligible (exercise 19 (h)). Show that the mapping \(\mathfrak{p} \mapsto \mathfrak{p}_{\mathcal{F}}\) is a bijection of the set of ideals \(\mathfrak{p}\) of \(\mathbf{A}\) not belonging to \(\mathcal{F}\) onto the set of prime ideals of \(\mathbf{A}_{\mathcal{F}}\) not belonging to \(\mathcal{F}'\) . (Show first that, if \(\mathbf{A}\) is an integral domain, so is \(\Lambda_{\mathcal{F}}\) , making the same note as in (a); then use the exact sequence of Exercise 19 (d) to show that \(\mathfrak{p}_{\mathcal{F}}\) is prime. If \(\mathfrak{p}'\) is a prime ideal of \(A_{\mathcal{F}}\) not belonging to \(\mathcal{F}'\) and \(\mathfrak{p} = j_{\mathrm{A}}^{-1}(\mathfrak{p}')\) , recall that \(\mathfrak{p}' = \mathfrak{p}_{\mathcal{F}}\) .
\item
  Let B be an A-algebra and G the set of left ideals l of B such that \(B_{s}/!\) is F-negligible. Show that G is the set of left ideals of B containing an ideal of the form B.m, where \(m \in F\) ; deduce that G is topologizing and idempotent and that, for every B-module N, FN = GN and \(N_{F}\) is canonically identified with \(N_{G}\) .
\item
  Let S be a multiplicative subset of A and F the set of ideals of A meeting S. Show that F is topologizing and idempotent and that, for every A-module M, \(M_{(F)}\) and \(M_{F}\) are canonically identified with \(S^{-1}M\) (note that every ideal of F contains a principal ideal As, where \(s \in S\) ).
\end{enumerate}

¶ 22. Let A be a ring (not necessarily commutative) and S a multiplicative subset of A.

\begin{enumerate}
\def\labelenumi{(\alph{enumi})}
\item
  Let \(\mathcal{F}\) be the set of left ideals \(\mathfrak{l}\) of \(\mathbf{A}\) such that, for all \(a\in \mathbf{A}\) , there exists \(s\in \mathbf{S}\) such that \(sa\in \mathfrak{l}\) ; this implies in particular that \(\mathfrak{l}\cap \mathbf{S}\neq \varnothing\) . Show that \(\mathcal{F}\) is topologizing and idempotent.
\item
  Let B be a ring and \(\phi: A \to B\) a ring homomorphism. The ordered pair \((B, \phi)\) is called a ring of left fractions of \(A\) with denominators in \(S\) if it satisfies the following conditions:
\end{enumerate}

\begin{enumerate}
\def\labelenumi{(\Roman{enumi})}
\item
  If \(\phi(a) = 0\) , there exists \(s \in S\) such that \(sa = 0\) .
\item
  If \(s \in S\) , \(\phi(s)\) is invertible in \(B\) .
\item
  Every element of B is of the form \((\phi(s))^{-1}\phi(a)\) , where \(a \in \mathbf{A}\) .
\end{enumerate}

Show that the following properties are equivalent:

(α) The ring A possesses a ring of left fractions with denominators in S.

(β) The following conditions hold:

\((\beta_{1})\) For all \(s \in S\) , \(a \in A\) , there exists \(t \in S\) and \(b \in A\) such that \(ta = bs\) .

\((\beta_{2})\) If \(a \in \mathbf{A}\) and \(s \in \mathbf{S}\) satisfy \(as = 0\) , there exists \(t \in \mathbf{S}\) such that \(ta = 0\) .

(γ) The canonical images of the elements of S in \(A_{F}\) are invertible.

Moreover these properties imply the following:

(δ) The principal ideals As, where \(s \in S\) , belong to F and every ideal \(m \in F\) contains one of these principal ideals.

(ε) The left annihilator of every \(s \in S\) is contained in \(\mathcal{FA}_{s}\) .

(ζ) For every exact sequence \(0 \to \mathbf{M} \xrightarrow{h} \mathbf{N} \xrightarrow{p} \mathbf{P} \to 0\) of A-modules, the sequence \(0 \longrightarrow \mathbf{M}_{\mathcal{F}} \xrightarrow{h_{\mathcal{F}}} \mathbf{N}_{\mathcal{F}} \xrightarrow{p_{\mathcal{F}}} \mathbf{P}_{\mathcal{F}} \longrightarrow 0\) is exact.

(Show that \((\alpha) \Rightarrow (\beta) \Rightarrow (\gamma) \Rightarrow (\alpha)\) . To see that \((\gamma)\) implies \((\alpha)\) , note that, for every A-module M and every \(x \in \mathcal{FM}\) , there exists \(s \in S\) such that \(sx = 0\) . To see that \((\beta)\) implies \((\gamma)\) , prove first that \((\beta)\) implies \((\delta)\) , \((\varepsilon)\) and \((\zeta)\) . To show that \((\beta)\) implies \((\zeta)\) , establish, using \((\delta)\) , that, if \(m \in \mathcal{F}\) and

\[
u \in \operatorname{Hom} _ {\mathbf {A}} (\mathfrak {m}, \mathrm{P} / \mathscr {F P}),
\]

then there exists \(s \in S\) such that \(As \subset m\) and \(v \in \operatorname{Hom}_{\mathbf{A}}(As, N/\mathcal{F}N)\) such that the diagram

\[
\begin{array}{c} \text {As} \xrightarrow {v} \text {N / FN} \\ i \Big \downarrow \qquad \qquad \qquad \qquad \qquad \qquad \qquad \qquad \Big \downarrow p ^ {\prime} \\ \mathfrak {m} \xrightarrow [ u ]{} \text {P / FP} \end{array}
\]

is commutative, i being the canonical injection and \(p'\) the mapping derived from p by taking quotients. Consider finally, for all \(s \in S\) , the exact sequence

\[
0 \longrightarrow \mathrm{N} \longrightarrow \mathrm{A} \xrightarrow {\mu_ {s}} \mathrm{A} \longrightarrow \mathrm{A} / \mathrm{As} \longrightarrow 0
\]

where \(\mu_s\) is the homothety \(\lambda \mapsto \lambda s\) ; use \((\varepsilon)\) , \((\zeta)\) and Exercise 4 of Algebra, Chapter I, § 8.)

\begin{enumerate}
\def\labelenumi{(\alph{enumi})}
\setcounter{enumi}{2}
\item
  Deduce from (b) that, if A possesses a ring of left fractions \((\mathbf{B}, \phi)\) with denominators in S, this ring has the following universal property: for every ring homomorphism \(f: A \to C\) such that \(f(s)\) is invertible in C for all \(s \in S\) , there exists a unique homomorphism \(g: B \to C\) such that \(f = g \circ \phi\) . In particular B is canonically isomorphic to \(A_{F}\) .
\item
  The notion of a ring of right fractions of A with denominators in S is similarly defined. Deduce from (c) that, if there exist both a ring of left fractions and a ring of right fractions of A, these two rings are isomorphic.
\item
  Let E be a vector space over a field K, \((e_{t})_{t\in I}\) a basis of E and T the tensor algebra of E. Let J be a non-empty subset of I distinct from I and S the multiplicative subset of T generated by the \(e_{t}\) such that \(t\in J\) . Show that T admits no ring of left or right fractions with denominators in S.
\item
  With the notation of (e) suppose that I = N and consider the quotient ring A of T by the two-sided ideal generated by the elements \(e_{0}^{2}\) and \(e_{0}e_{i} - e_{i+1}e_{0}\) , for all i \textgreater{} 0, and \(e_{i}e_{j} - e_{j}e_{i}\) , for i \textgreater{} 0, j \textgreater{} 0. Let \(S'\) be the multiplicative set generated by the canonical images of the \(e_{i}\) in A, for i \textgreater{} 0; show that A admits a ring of left fractions, but not a ring of right fractions, with denominators in \(S'\) .
\item
  Suppose that A admits a ring of left fractions with denominators in S. In the set \(S \times A\) , let R be the relation between \((s, x)\) and \((t, y)\) : ``there exists \(r \in S\) , \(u \in A\) , \(v \in A\) such that r = us = vt and ux = vy''. Show that R is an equivalence relation and define on the set \(B = (S \times A)/R\) a ring structure such that the ordered pair consisting of S and the mapping \(\phi: A \to B\) , which maps \(x \in A\) to the class of \((1, x)\) , is a ring of left fractions of A with denominators in S (use the fact that, for \(s \in S\) and \(t \in S\) , there exists \(u \in A\) and \(v \in A\) such that r = us = tv belongs to S). Consider the case \(S = A - \{0\}\) (cf.~Algebra, Chapter I, §9, Exercise 8).
\end{enumerate}

\begin{enumerate}
\def\labelenumi{\arabic{enumi}.}
\setcounter{enumi}{22}
\tightlist
\item
  \begin{enumerate}
  \def\labelenumii{(\alph{enumii})}
  \tightlist
  \item
    Let A be a left Noetherian ring with no divisor of 0. Show that A possesses a field of left fractions (use Exercise 14 (a) to show that condition \((\beta_{1})\) of Exercise 22 (b) is fulfilled for \(S = A - \{0\}\) ).
  \end{enumerate}
\end{enumerate}

\begin{enumerate}
\def\labelenumi{(\alph{enumi})}
\setcounter{enumi}{1}
\tightlist
\item
  Let A be a left and right Noetherian ring with no divisor of 0. Show that every finitely generated left A-module M, each of whose elements \(\neq 0\) is free, is a submodule of a free A-module. (Embed M canonically in \(\mathbf{K} \otimes_{\mathbf{A}} \mathbf{M}\) , where K is the field of fractions (left or right) of A considered as a right A-module; if \((x_{i})\) is a basis of the left vector K-space \(\mathbf{K} \otimes_{\mathbf{A}} \mathbf{M}\) , show that there exists \(a \neq 0\) in A such that M is contained in the sub-A-module of \(\mathbf{K} \otimes_{\mathbf{A}} \mathbf{M}\) generated by the \(a^{-1}x_{i}\) .)
\end{enumerate}

¶ 24. Let A be a ring, commutative or not, and F the set of left ideals 1 of A such that \(l \cap m \neq 0\) for every left ideal \(m \neq 0\) of A.

In order that \(l \in F\) , it is necessary and sufficient that, for every element \(a \neq 0\) of A, there exist \(b \in A\) such that \(ba \neq 0\) and \(ba \in l\) . Every ideal l of F contains the left socle of A (Algebra, Chapter VIII, §5, Exercise 9) and, if A is Artinian, F consists of the ideals containing the left socle of A.

\begin{enumerate}
\def\labelenumi{(\alph{enumi})}
\tightlist
\item
  Show that \(\mathcal{F}\) is a topologizing set (Exercise 16). Deduce that the set \(\mathfrak{a}\) of \(a \in \mathbf{A}\) such that \(\mathfrak{l}a = 0\) for some \(\mathfrak{l} \in \mathcal{F}\) (union of the right annihilators of the \(\mathfrak{l} \in \mathcal{F}\) ) is a two-sided ideal of \(\mathbf{A}\) , which contains no idempotent. If \(\mathbf{A}\) is left Noetherian, show that \(\mathfrak{a}\) is nilpotent (prove first that every element \(a \in \mathfrak{a}\) is nilpotent, by considering the left annihilators of the \(a^n\) ; then use Exercise 26 (b) of Algebra, Chapter VIII, § 6).
\end{enumerate}

Let K be a commutative field and B the ring of polynomials with coefficients in K in indeterminates Y and \(X_{i}\) ( \(i \geqslant 1\) ). Let b be the ideal of B generated by the \(X_{i}X_{j}\) ( \(i \neq j\) ) and the \(Y^{i-1}X_{i}\) and let A be the quotient ring B/b; show that in A the ideal a defined above contains the class of Y, which is not nilpotent.

\begin{enumerate}
\def\labelenumi{(\alph{enumi})}
\setcounter{enumi}{1}
\item
  Show that, for every left ideal m of A, there exists a left ideal n of A such that \(m \cap n = 0\) and \(m + n \in F\) . (Take n to be an ideal which is maximal amongst the left ideals l such that \(m \cap l = 0\) .)
\item
  A ring A is called (left) neat if the two-sided ideal a defined in (a) is reduced to 0. If this is so, the set F is idempotent (let l, m, n be left ideals of A such that \(n \in F\) , \(l \subset n\) , \(n/l\) being F-negligible, and \(m \neq 0\) . If \(x \neq 0\) belongs to \(m \cap n\) , there exists \(r \in F\) such that \(r.x \subset l \cap m\) ).
\end{enumerate}

¶ 25. With the notation of Exercise 24, suppose that A is a left neat ring.

\begin{enumerate}
\def\labelenumi{(\alph{enumi})}
\item
  Show that the canonical mapping \(j_{\mathbf{A}}: \mathbf{A} \to \mathbf{A}_{\mathcal{F}}\) is injective, which allows us to identify A with a subgroup of \(\mathbf{A}_{\mathcal{F}}\) . Let \(\mathcal{F}'\) be the set of left ideals \(l'\) of \(\mathbf{A}_{\mathcal{F}}\) such that \(l'_{\mathcal{F}} = \mathbf{A}_{\mathcal{F}}\) ; show that \(\mathcal{F}'\) is the set of left ideals \(l'\) of \(\mathbf{A}_{\mathcal{F}}\) such that \(l' \cap m' \neq 0\) for every left ideal \(m' \neq 0\) of \(\mathbf{A}_{\mathcal{F}}\) . (Show first that, for every left ideal \(m' \neq 0\) of \(\mathbf{A}_{\mathcal{F}}\) , \(A \cap m' \neq 0\) ; note on the other hand that the ideals \(l' \in \mathcal{F}'\) are precisely the left ideals of \(\mathbf{A}_{\mathcal{F}}\) containing an ideal \(m \in \mathcal{F}\) .) Show that the ring \(\mathbf{A}_{\mathcal{F}}\) is left neat.
\item
  Show that \(\mathbf{A}_{\mathcal{F}}\) is an injective A-module, that is (Algebra, Chapter II, §2, Exercise 11) that, for every left ideal I of A and all
\end{enumerate}

\[
u \in \operatorname{Hom} _ {\mathbf {A}} (\mathfrak {l}, (\mathrm{A} _ {\mathcal {F}}) _ {s}),
\]

there exists \(a' \in \mathbf{A}_{\mathcal{F}}\) such that \(u(x) = xa'\) for all \(x \in \mathfrak{l}\) . (Use Exercise 24 (b) and Exercise 19 (c).) Deduce that every endomorphism of the A-module \(\mathbf{A}_{\mathcal{F}}\) is of form \(u: x \mapsto xa'\) where \(a' \in \mathbf{A}_{\mathcal{F}}\) (note that \(\operatorname{Hom}_{\mathbf{A}}(\mathbf{A}, \mathbf{A}_{\mathcal{F}}) = \operatorname{Hom}_{\mathbf{A}_{\mathcal{F}}}(\mathbf{A}_{\mathcal{F}}, \mathbf{A}_{\mathcal{F}})\) by Exercise 19 (e)); also \(\operatorname{Ker}(u) = (\operatorname{Ker}(u))_{\mathcal{F}}\) .

\begin{enumerate}
\def\labelenumi{(\alph{enumi})}
\setcounter{enumi}{2}
\item
  Show that, for every left ideal \(\mathfrak{l}\) of \(\mathbf{A}\) , \(\mathfrak{l}_{\mathcal{F}}\) is a direct factor of \(\mathbf{A}_{\mathcal{F}}\) . (Let \(\mathfrak{m}\) be a left ideal of \(\mathbf{A}\) such that \(\mathfrak{m} \cap \mathfrak{l} = 0\) and \(\mathfrak{m} + \mathfrak{l} \in \mathcal{F}\) ; extend the projection of \(\mathfrak{m} + \mathfrak{l}\) onto \(\mathfrak{m}\) , which is zero on \(\mathfrak{l}\) , to an endomorphism \(p\) of \(\mathbf{A}_{\mathcal{F}}\) ; show that \(p^2 = p\) and that \(\operatorname{Ker}(p) = \mathfrak{l}_{\mathcal{F}}\) .)
\item
  Show that \(A_{F}\) is an absolutely flat ring (Chapter I, §2, Exercise 17). (Reduce it to the case where \(A = A_{F}\) . If u is the endomorphism \(x \mapsto xa\) of the left A-module \(A_{s}\) , then \(\operatorname{Ker}(u) = (\operatorname{Ker}(u))_{\mathcal{F}}\) , and \(\operatorname{Ker}(u)\) therefore admits a complementary ideal m in \(A_{s}\) ; note that \(\operatorname{Im}(u)\) is isomorphic to m, hence injective (Algebra, Chapter II, §2, Exercise 11) and therefore a direct factor of \(A_{s}\) (loc. cit.). Finally, use Exercise 13 of Algebra, Chapter II, §1.)
\end{enumerate}

¶ 26. (a) Show that every Zorn ring A (Algebra, Chapter VIII, § 6, Exercise 13) with no nil ideal ≠0 is left neat; in particular, every absolutely flat ring and every primitive ring whose socle is not reduced to 0 is a left neat ring.

\begin{enumerate}
\def\labelenumi{(\alph{enumi})}
\setcounter{enumi}{1}
\item
  Let A be a left primitive ring whose socle is not reduced to 0; show that, if A is represented as a dense subring of a ring of endomorphisms \(\operatorname{End}_{\mathbb{D}}(\mathbf{V})\) of a vector space (Algebra, Chapter VIII, § 5, Exercise 10), the ring \(A_{F}\) corresponding to A (Exercise 25) is isomorphic to \(\operatorname{End}_{\mathbb{D}}(\mathbf{V})\) .
\item
  Show that every quasi-simple ring A (Algebra, Chapter VIII, § 5, Exercise 5) is left and right neat and that the corresponding ring \(\mathbf{A}_{\mathcal{F}}\) is quasi-simple. Deduce that, for every ring B not reduced to 0, there exists a homomorphism \(\phi: \mathrm{B} \to \mathrm{R}\) to a ring R which is quasi-simple, absolutely flat and such that \(\mathbb{R}_s\) is an injective R-module.
\item
  Let A be a left Noetherian ring with no nilpotent ideals. Show that A is left neat and that the corresponding ring \(A_{F}\) is a semi-simple ring (``Goldie's Theorem''; use Exercise 24 (a) and, on the other hand, use Algebra, Chapter VIII, § 2, Exercise 7, noting that in the absolutely flat ring \(A_{F}\) there cannot exist any family of idempotents which are pairwise orthogonal, unless A contained an infinite direct sum of left ideals \(\neq0\) ). Show that, if \(s\in A\) is not a right divisor of 0 in A, it is invertible in \(A_{F}\) and conversely; the set S of these elements is a multiplicative subset of A. Show that F is the set of left ideals of A which meet S (to see that the ideals of A meeting S belong to F, note that the elements of S are not left divisors of 0 in A; to see that every \(l\in F\) contains an element of S, represent each simple component of \(A_{F}\) as an endomorphism ring of an n-dimensional vector space E and define by induction n elements \(u_{1},\ldots,u_{n}\) such that \(u_{i}u_{j}=0\) for j\textless i and \(u_{i}^{2}\neq0\) , the \(u_{i}\) being of rank 1; finally show that the kernel of \(s=u_{1}+\cdots+u_{n}\) is zero). Deduce that \((\mathbf{A}_{\mathcal{F}},j_{\mathrm{A}})\) is a ring of left fractions of A with denominators in S (Exercise 22).
\end{enumerate}

\begin{enumerate}
\def\labelenumi{\arabic{enumi}.}
\setcounter{enumi}{26}
\tightlist
\item
  \begin{enumerate}
  \def\labelenumii{(\alph{enumii})}
  \tightlist
  \item
    Let A be a ring and M a finitely generated faithful A-module. Show that, if a, b are two ideals of A such that aM = bM, the radicals of a and b are equal (use no. 2, Corollary 2 to Proposition 4).
  \end{enumerate}
\end{enumerate}

\begin{enumerate}
\def\labelenumi{(\alph{enumi})}
\setcounter{enumi}{1}
\item
  Deduce from (a) that, if in an integral domain A p is a finitely generated prime ideal, then \(p^{m} \neq p^{n}\) for \(m \neq n\) .
\item
  Let K be a field and A the polynomial ring K{[}X, Y{]} in two indeterminates. Let a be the ideal \(AX^{4} + AX^{3}Y + AXY^{3} + AY^{4}\) , b the ideal \(a + AX^{2}Y^{2}\) of A; show that \(b \neq a\) and \(b^{2} = a^{2}\) .
\end{enumerate}
